\chapter{YouTube}

\section{Scope and Goals}

As a kid, I had always wanted a YouTube Channel. Now I will be working towards that, on a bit more pratical level. For the purposes of learning more personally, leveling up my skill as a

\section{Getting Started}

\subsection{Ideas}

At some point I want a dive into Kerr Black holes but I am not ready for that.

\begin{itemize}
  \item Numerical Hamiltonian(done)
  \item A Physical Look Into Hamiltonian Mechanics
  \item Hidden Assumptions in Modern Physics
  \item Journey through General Relativity
  \item Positive Mass Warp Drive
  \item Axioms of Hamiltonian Mechanics
  \item Understanding of Metriplectic Dynamics
  \item Gravitational Waves
  \item MagnetoHydrodynamics Propulsion
  \item MagnetoHydrodynamics in a nebulua
  \item Deeper Numerical Methods
  \item Scientific Programme
  \item Destroying a Universe with a punch(done)
  \item Best way to simulate a solar system(done)
  \item Metriplectic simulation
  \item In-depth operator splitting(done*)
  \item Deriving Plasma kinetic equations
  \item Simple GR solver
  \item ADM
  \item Symplectic geometry: A physicist's perspective
  \item DG
  \item Deriving Gravitational Wave Memory
\end{itemize}


\section{Introduction to Symplectic Integrators}

\subsection{General overview of direction}
\begin{itemize}
\item Introduction to me
\item Motivation
\item Hamiltonian Mechanics introduction
\item Geometric perspective
\item Symplectic Euler compared to traditional Euler
\item Leapfrog/Verlet
\item Practicule guidence
\end{itemize}

\subsection{Opening}

Before I begin this video; I am not an expert. Must of my understanding comes from Problems In Hamiltonian Dynamics by (.....). This is more for my own education than others, for a better understanding of the concepts and development of my teaching abilities. I highly recommened continue to learn far beyond this simple introduction I will lay out.
(picture of the book)
\subsection{Motivation}
Symplectic Numerical methods are focus inheritily on the structure and methods of Hamiltonian mechanics. Thus, to understand them, one must understand why. Why it is used at all. At first glance, it seems quite contrived. Hamiltonian Mechanics is quite mathematically abstract and complicated. Beyond that, expanding it in differing ways such as dissipative or entropic spaces requires creative solutions. Thus, begs the question, why is it such a developed field?

(put really complex hamiltonian)

The short answer is that Hamiltonian Mechanics geometry enforces energy conservation in a way others formulisms do not. The long answer... is well this video.

(put graph showing energy conservation)
\subsection{Hamiltonian Mechanics Review}
First, a quick review on Hamiltonian Mechanics. If you are comfortable feel free to skip ahead.

A hamiltonian system is defined by the function H(q,p,t) - the hamiltonian - which repersents something analogus to total energy. Not exactly though, hopefully in a future video we will get to that, but for not thinking of it in energy works great for this purpose.
\[
H(q,p,t) = T + V
\]
Here q repersents generalized position, which is the minimum amount of set parameters that fully spesify where a system is in configuration space. Where configuration space is a space whose points repersent all possible configurations. Meaning that generalized position is not confined to distance, but also things like angles, diminsionless quantaties, or what ever other convenient quantities that can be used.

The q, repersents generalized momenta. Which holds the same rules as position does, but repersenting momenta instead of position.

The evolution can thus be derived to be
\[
\frac{dq}{dt} = \frac{\partial H}{\partial p}
\]
\[
\frac{dp}{dt} = -\frac{\partial H}{\partial p}
\]
Let us look at a concrete example: a simple pendulum

(picture of pendulumn)

\[
H(q,p) = \frac{p^2}{2m} + \frac{1}{2} kq^2
\]

Such, one can take the partial derivatives to find:
\[
\frac{dq}{dt} = \frac{p}{m}
\]
\[
\frac{dp}{dt} = -kq
\]

In phase-space, the space of all generalized pairs, the system traces out a tragectory. For harmonic ossilators these are ellipses. Each ellipse corresponds to a different energy level.
(show phase-space)

Here are some key points.

- energy remains constant along the tragectories. As tragectories are defined by conserved energy

- the phase-space volume is conserved. This is essentially
(shade in one region)

-finally, the most importaint for this, the manifold retains a symplectic structure. Which is defined at the 2-form $ \omega = dq \wedge dp$. This will be disgussed further shortly.

\subsection{Why standard models fail}
To understand the symplectic methods, one must find their motivation. Why traditional methods fail in certain senerios, no matter the type. To understand this, we will look at the simplest integrator. Forward Euler

\[
q_{n+1} = q_n + h * \frac{\partial H}{\partial q}_n
\]
\[
p_{n+1}=1_n +h * \frac{\partial H}{\partial q}_n
\]
Where h is the time step.

Let us apply this to our harmonic ossilater and see what happens... Tragectory spirals outwards! This is simply because traditional methods aren't made to conserve energy, they are meant to model physics. The main reason is that tiny errors add up without the ability to contain them. Also, traditional methods add unphysical correction terms that over time lead to unphysical effecs. Being dissipate in certain areas and adding energy in others. Adding artificial energy to the system in order to have temporary  Now, obviously higher order corrections help with these problems and making the models more physical helps, but these problems are fundemental. Adding higher order corrections simply increases the computational workload without ever fully resolving the problem.

Even sofisticated methods like RK4 have this problem. It is accurate to the 4th order, but that means that the error per step is $O(h^5)$. This is tiny, but adds up. For something like a celestial simulation over millions or even billions of years, this creates distorted simulations.

(Log-log plot showing error growth)

\subsection{Symplectic Integrators}
(illistration of symplectic maps)

Symplectic integrators fix this problem by peserving the symplectic structure I mentioned earlier.

\[
\omega = dq \wedge dp
\]

The symplectic form is a mathematical object that mesures oriented areas in phase-space. For a simple system with one degree of freedom in the position and momentum each one gains this equation.

This is a 2-form, when integrated over any region of phase space, you get an oriented area.
(physical deformation)

As the hamiltonian flow evolves, the region over phase-space gets stretched, rotated, and deformed. But, the symplectic form remains perserved.

From this basic axiom, one can derive something called the Shadowing Theorem, that says:

A tragectory computed by symplectic integrator, while not exactly true, stays close to the tragecory of a slightly modified Hamiltonain.

(visual: compared integration to true)

Meaning, standard integrators essentially create a new equation that is close to the origional.

Symplectic integrators create another hamiltonian that is exactly equal to the origional hamiltonain plus the order of the error.

(energy conservation plot)

Such, the energy does not increase monotonically, it ocilates within a bounded region. Over exponentially long time, the error grows logorithmicly rather than linearly.

\subsection{Symplectic Euler's method}

To compare, we will once again take the simpliest possible case.

Take the traditional Forward Euler
\[
q_{n+1} = q_n + h * \frac{\partial H}{\partial q}_n
\]
\[
p_{n+1}=1_n +h * \frac{\partial H}{\partial q}_n
\]

Now turn it into what is called a backward Euler
\[
q_{n+1} = q_n + h * \frac{\partial H}{\partial q}_n
\]
\[
p_{n+1}=1_n +h * \frac{\partial H}{\partial q}_n
\]
Notice what has changed, we updated p using the old position but the momentum by using the new momentum. One thing to note, backwards or implicit integration is commonly used outside of symplectic integration, it is simply a general requirments for symplectic integration. There are some special cases where explcit numerical methods satisfy the requirments of a symplectic integrator. Another requirment being that one using a Hamiltonian or Hamiltonian like formulism such as Numbu mechanics.

Now let us see why it is symplectic:
\[
H(q,p) = T(p) + V(q)
\]

where T is kinetic energy and V potential.

The symplecti Euler step can be decomposed into pure position-step, which is the flow of the potential. The a pure momentum step, which is exactly the flow of kinetic energy. The composition of these is the exact flow of the Hamiltonian, which is by definition symplectic.

The exact requirments are quite involved and mathematically difficult to derive, so an independent video on the topic is reqiured. For the purposes of this video the general rules described and the types of integrators are suffient.

Let us implement this for the harmonic oscillator.

In this, the orbit remains stable. Energy oscilates but doesn't drift.

One can also make the equation implicit by updating q before p.

\subsection{The LeapFrog/Verlet Method}

The most commonly used sympletic integrator is the leapfrog method. It is second order accurate.

The core idea: Evaluate positions and momenta at staggered time points.

(equation)

(visual)

Think of it as half-step momentum backward. Full step position forward, using the half step. Then half-step momentum backward again.

This is symmetric: if you run it backward, you get the exact same algorithm. This time-reversibillity is what gives second order accuracy.

There is an equivlent position formulism that is often more convenient.
(equation)

Where v is velocity and a is acceleration. Note, that while it uses variables commonly used in Newtonian dynamics, it is still fundementally hamiltonian.
(visual)
One can compare performence after 1000 orbits. For a timestep of 0.01 Forward euler has energy error of 15.3\%, symplectic euler 0.8\%, and leapfrog 0.002\%.

As you can see much smaller, as the error scales by the square of the time step. Allowing for larger time steps for the similar accurecy.

A derivation of this equation is outside of the scope of this video but is an interesting idea that may be develed into in the future.

\subsection{Splitting Method}

To build higher order symplectic integrators, one must learn operator splitting.

The key idea: if you can write the Hamiltonian as a sum
\[
H=H_A + H_B
\]

and you know how to integrate each piece exactly, you can build approximate integrators by composing these exact functions of these two.

where $\phi_A(h)$ is the exact flow generated by $H_A$ over time h.
and $\phi_A$ is the exact flow over $H_B$.

Then the true hamiltonian flow is:



%\[
%\phi_H(h) = \exp\!\left(h\,\{H,\cdot\}\right)
%\]




where $\{\cdot, \cdot \}$ denotes the Poisson Bracket.

The simpliest approximation is obtained by composing the two exact flows, this is known as the Lie-Totter splitting, this method is first otrder accurate. Higher order methods can be built using splitting, such as the 4th-order Forest-Ruth:

\[
\psi_4(h)=\psi_2(\theta h) \circ \psi_2((1-2 \theta)h) \circ \psi_2(\theta h)
\]

A derivation of this is outside the scope of this video but is something we will get into.

\subsection{Beyond the Basics}
Symplectic integration goes far beyond the introduction given in this video.

\begin{itemize}
\item Dynamics on a manifold
\item veriational integration
\item Backward error analysis
\item Adaptive symplectic methods
  \item Methods for Hamiltonian adjacent frameworks(ie. Nambu, metriplectic, etc)
\end{itemize}

Some of these I would like to explore in future videos but many of them are currently beyond my expertice.

\section{Physical Interpretations of Hamiltonian Dynamics and Byond}

\subsection{Planning}

\begin{enumerate}
  \item Opening
  \item What is the Hamiltonian
  \item What physically is phase space
  \item Dissipation and the Non-Hamiltonian World
  \item Beyond Symplectic: Metriplectic and Generic
  \item Beyond Symplectic: Contact Geometry and Geometrization of Irreversibility
  \item Beyond Symplectic: Poisson Manifolds, Casimirs, and Nambu
  \item Structural Realism
  \item Open Problems
\end{enumerate}

\subsection{Opening}

Often in higher level mathematical physics, you will here many phrases thrown around with generally little to no explanation. ''the hamiltonian is analougus to total energy but isn't literally'', ''Hamiltonian mechanics forbids dissipation'', and many more.

These complex ideas are pushed aside for latter, well if you are like me, you think that later is now, so that is the purpose of this video.

This video will require a good intuitive and understanding of Hamiltonian mechanics and the basic structural underpinnings. Also, be warned, this is a fun project, mistakes are entirely possible.

\subsection{What is the Hamiltonian?}

In most courses, you will be introduced to the Hamiltonian functions through this equation
\[
H = V + T
\]
where the Hamiltonian is the addition of kinetic and potential energy, and is thus the total energy.

A few years later you may here terms like generator of the flow, with only the explanation that H is a function that generates the correct observables, what ever those are.

So, how do these definitions fit together, and what do they even mean.

Essentially, the idea is that H uses Noether's theorem's time translation to make it so that the Hamiltonian is always equal to:

\[
\pdv{A}{t} = \left\{A,H  \right\}
  \]
  where H is the observable. An observable is a pretty self explanatory concept. Something that can be physically observed; momentum, energy, velocity, entropy, or even non-physics measurements like personality. Though classical Hamiltonians have trouble with entropy as we will get into.

  As you can figure out with some careful playing; the Hamiltonian is only equal in a special case; in classical systems where the observables are position and momenta. In other cases, it no longer is that. It could be something else, or simply a number with no traditional physical analogy.



\subsection{What is Phase Space?}
Now, to actually use the Hamiltonian, knows we need phase space. So, what physically is that.

Well, a good analogy is imagine the universe is frozen in time, what information do you need to predict what happens next?


Now that is just an analogy, if you are watching this video I suspect you want something more than that.

Well, we must first start with a state, which is really just a property, the properties needed for physics. Newton defined these properties through forces, mass, position, time, and their respective derivatives to arbitrary orders.

Hamilton, defined it through configuration momenta and position. It really can be anything but as we will motivate later in most physical systems these are the generally the most useful. Then, he found you are able to for every possible physical state, there could be one point in this phase space of position and momenta that corresponds to it. Such as a 2D body can be ploted in 4D space through
\[
(x,y,p_x,p_y)
\]

There, objects no longer viewed as moving through space but rather tragectories in phase space. Generally this view is considered to be a mathematically significant but metaphysically trivial change, but some people disagree, which is something that will be covered later in the video.

\subsection{Is Momentum more Fundamental than Velocity}

The question of which is more fundemental; momentum or velocity seems quite arbitrary, and in many ways it is, but we will look into it regardless.

In traditional Newtonian Mechanics; momentum is traditionally derived by velocity, but the inverse is just as true.

In hamiltonian mechanics, this isn't the case. Momentum is more than just mass and velocity, it also includes think like electric fields or even other fundemental forces like gravity in some cases. In quantum mechanics this connection becomes even more muddled for obvious reasons.

Now to really answer the question. I am sad to say that it depends on what you are looking for. Velocity is much more intuitive, but momentum is a generator of translations and naturally connected with symplectic structures. Likely not a very satisfying answer, but the only one I have.


\subsection{Hamiltonian Mechanics and Teleology}


\subsection{Noether's Theorem}


\subsection{Canonical Transformations}

\subsection{The Question of Dualism}

\subsection{Hamiltonian Branch offs}

\subsection{Hamiltonian and Dissipation}

\subsubsection{Metriplectic and Others}

\subsection{Structural Realism}
%%%%%%%%%%%%%%%%%%%%%%%%%%%%%%%%%%%%%%%%%%%%%%%%%%%%%%%%%%%
\section{Variational Integrator for RMHD}

\subsection{Planning}
\begin{enumerate}
  \item Opening
  \item Introduction
  \item RMHD
  \item Formal Lagrangian
  \item Conservation Laws
  \item Electron Inertia
  \item Symmmetrisation
  \item Discrete Action Principle
  \item Discrete Conservation Laws
  \item Numerical Methods
  \item Validation
  \item Random Testing
\end{enumerate}

I am thinking for random testing either a video of one of the earlier simulations or do something random.


\subsection{Opening}

Before I begin, as always, I am not an expert. Most of my information comes from the paper ``Variational Integrators for Reduced Magnetohydrodynamics'', by Michael Kraus, Emanuele Tassi, and Danueka Grasso. For more information their paper will be in the description. Anyways, on to the video.

\subsection{Introduction}

For all purposes, the subject of this video is the creation of a numerical method that takes a simplified model of magnetized plasma and preserves the conservation laws with greater accuracy than most other methods. The tool used to do so is called a variational integrator. That descretizes the action, as in the action in Lagrandian Mechanics, rather than the equations of motion.
\subsection{Variational Integrators}
Variational integrators work in the inverse of traditional methods. In traditional methods you start with the model
\[
F = kx
\]
then define the equations of motion
\[
\ddot{x} + \frac{k}{m}x = 0
  \]

  From here one descretizes the equation of motion.

  For variational derivatives one takes the action:
\[
S[q] = \int_{t_0}^{t_1} L(q, \dot{q}) dt
  \]

  The nature selects where
\[
  \delta S = 0
  \]

  From here derivation of things such as
  \begin{itemize}
    \item Euler-Lagrange
    \item Hamiltonian
    \item Conservation laws
  \end{itemize}
  From here instead of discretizing the differential equation, one discretizes the lagrangian
\[
S_d = \sum_{k=1}^{N-1} L_d(q_k, \dot{q_{k+1}})
  \]
  Where
\[
L_d \approx \int_{t_k}^{t_{k+1}}L dt
  \]

  From here, varing across the sequence rather than a smooth curve gives you the Discrete Euler-Lagrange Equation. Here the equation is already an algorithm, no additional methods needed(though we will use some in this video)
\subsection{RMHD}

To do this, we must define what RMHD is, to do that we must first define MHD. MHD treates plasma as a single conductive fluid coupled with a magnetic field. While it is a massive simplification of ``Billions of charged particles'' it is still quite accurate and obviously much less computationally expensive.

RMHD, further simplifies, for a spesific, useful physical situation. In plasma with a strong backgroud magnetic field, like in a tokamak, where the field mostly points one way (call it "toroidal") and everything interesting happens in the plane perpendicular to it. Under those conditions, the fast physics along the strong field direction basically stops mattering, and you're left with an effectively two-dimensional problem.

With only four fields remaining as useful:
\begin{itemize}
  \item $\phi$, the stream function, describes how the fluid flows
  \item $\omega$, the vorticity, how the flow swirls
  \item $\psi$, the magnetic potential
  \item \textbf{j}, the current density, how the magnetic field swirls
\end{itemize}

The equation then simplifies to:
\[
  \omega_t \pb{\phi}{\omega} + \pb{j}{\psi} = 0, \quad - \Delta \phi = \omega
\]

\[
\psi_t + \pb{\phi}{\psi}= 0, \quad - \Delta \psi = j
  \]
  With the poisson bracket being simply two dimensional so that

\[
\pb{\phi}{\omega} = \phi_x \omega_y - \phi_y \omega_x
  \]

  Here, is can be seen that RMHD has Casimir Invariants, quantaties that are conserved, not from traditional physical symmetries but from the structure of the Lie algebras. There are families with infinite many subsets, for the pusposes of thisl we will focus on Magnetic helicity(how tangled the magnetic fields are) and Cross helicity(how aligned the flow is with the magnetic field)


  Now, for this paper to work, rather than use Hamiltonian structure, it uses Lie-Poisson, Hamilton's weird cousin. It still conserves energy but standard symplectic methods simply do not work.

  This is because fluid mechanics observables do not naturally describe motion along a phase space, but live on a Lie Algabra with symmetries within it. So the Poisson operator changes with the state rather than is a constant symplectic structure.

\subsection{Formal Lagrandian}

From here comes a snag, variational integrators need a Lagrandian, an action to discreize. But, for RMHD, there is no obvious Eulerian-Lagrandian. So what does one do?

They make one up. There is a trick called a formal lagrandian. Essentialy, take every equation equation in your system and multiply is by a made up variable called an auxiliary variable that has no physical meaning. Then add them up
\[
L = \xi\left(\omega_t + \{\phi,\omega\} + \{j,\psi\}\right)
    + \mu\left(\omega + \Delta\phi\right)
    + \chi\left(\psi_t + \{\phi,\psi\}\right)
    + \zeta\left(j + \Delta\psi\right)
 \]

 Here four variables $ \xi, \mu, \chi, \zeta $, are all ghost variables that only exist to make the arithmatic work.

This helps by giving stationary variables to allow for one to vary along the physical variables using Hamilton's Principle leading to a brand new ``adjoint equation.''

Leading to the Euler-Lagrandian:
%\[
%\frac{\partial L}{\partial \varphi_a}\!\left(\varphi, \varphi_t,  \varphi_x, \varphi_y\right) %
%- %
%\frac{\partial}{\partial x^\mu} %
%\left(\frac{\partial L}{\partial (\partial_\mu \varphi_a)}\!\left(\varphi, \varphi_t, \varphi_x, \varphi_y\right)\right) %
%= 0 %
%\quad\text{for all } a. %

 % \]

  Where $\varphi$ denotes all of the variables. Many of the auxillary variables can be found to be to be functions of the origional variables. So that $\varphi$ can be defined as
\[
\varphi = (\omega, \psi, \phi, j, \psi \omega, 0,0)
  \]

\subsection{Conservation Laws}

Now, the whole puspose of the creation of the Lagrandian was to apply Noether's theorem. We restrict our attention to conservation laws which are generated by vertical
transformations, that is transformations which only affect the field variables but leave the base space invariant. From here with some complex math I won't get into, it is found that the symmetries are
\begin{itemize}
  \item Magnetic Helicity
  \item $L^2$ norm of $\psi$
  \item Cross Helicity
  \item Energy
\end{itemize}

\subsection{Electron Inertia}
This current RMHD, has a defect. It does not include magnetic reconnection, field line snapping and reconnecting. Which tend to be one of the most importaint studies in plasma physics. So, we will add a small correction, electron inertia. Real electrons have mass, just a small one, ones you treat them as such a new term develops in Ohm's law, one that allows reconnection.
Thus the new model changes $\psi$ to $\overline{\psi}= \psi+d^2_e j$ to give the electron canonical momentum where $d_e$ denotes the electron skin depth.

That's it, that is the whole mortification. From here, small bit of derivation shows that the Hamiltonian stays the same, so does the Casimir invariants and the conservation laws, only expressed in for of $\overline{\psi}$ rather than $\psi$.

\subsection{Symmetrisation}

Before we can put anything into our computer, there is some tedious house keeping needed. Our formal Lagrandian has a second derivative, which is highly challenging to discretize directly and even more costly to compute. To do this we use symmetrisation, essentially integration by parts but for functional rather than function's. We do this on the terms that produce the Laplancians, and get an equivalent Euler-Lagrange equation that only depends on the first order derivative of the fields.


Further care can be taken when discretising the Poisson Bracket. It can be seen through integration by parts when assuming appropriate boundary conditions that the cyclic permutations of the functions in the integrandare all identical

Insted of selecting on of those forms, a convex combination can be considered.
\[
\int \xi \pb{\phi}{\omega} dx dy = \int \left[\alpha \xi \pb{\phi}{\omega} + \beta \phi \pb{\omega}{\xi} + \gamma \omega \pb{\xi}{\phi}   \right]
  \]
  Where $\alpha = \beta = \gamma = \frac{1}{3}$ to conserve symmetry. From here we find the motified Lagrangian

\begin{align}
L^{\prime\prime\prime}(\varphi, \varphi_t, \varphi_x, \varphi_y)
&= \frac{1}{2}\,\varphi_t^{T}\Lambda\,\varphi \\
&\quad + \frac{1}{3}\left(
    \xi\{\phi,\omega\} + \phi\{\omega,\xi\} + \omega\{\xi,\phi\}
\right) \\
&\quad + \frac{1}{3}\left(
    \xi\{j,\psi\} + j\{\psi,\xi\} + \psi\{\xi,j\}
\right) \\
&\quad + \frac{1}{3}\left(
    \chi\{\phi,\psi\} + \phi\{\psi,\chi\} + \psi\{\chi,\phi\}
\right) \\
&\quad + \mu\,\omega - \nabla\mu\cdot\nabla\phi
      + \zeta\,j - \nabla\zeta\cdot\nabla\psi .
\end{align}


\subsection{Discrete Action Principle}

As mentioned earlier we will discretize the action rather than the equation of motion, for simplicity we will use Veselove-type finite difference to split into two space and one time. But the extention to more elaborate approaches is fairly straight forwards from here.

Here we can quite easily find that
%\begin{align}
 % \int_{x_i}^{x_{i+1}} dx \int_{y_i}^{y_{i+1}}dy\,
  %L(\varphi, \varphi_t, \varphi_x, \varphi_y)
  %&\approx
  %L_h\!\left(
   %   \varphi_{i,j}(t), \varphi_{i+1,j}(t), \varphi_{i+1,j+1}(t), \varphi_{i,j+1}(t), \right.\\
 % &\qquad\quad\left.
  %    \dot{\varphi}_{i,j}(t), \dot{\varphi}_{i+1,j}(t),
   %   \dot{\varphi}_{i+1,j+1}(t), \dot{\varphi}_{i,j+1}(t)
  %\right).
%\end{align}

A large deal of algebra allows us to get to the point of employing the Hamilton-Pontryagin Principle. Which states when transfering from Lagrandian to Hamiltonian we should not assume velcity and momentum's relation but derive it after the fact of treating them as independent.

The using what is called the Arakawa's discretization for the Poisson bracket, which is given by
\begin{align}
A_{i,j}(\phi,\omega)
&= \frac{1}{3} \Big(
    A^{++}_{i,j}(\phi,\omega)
  + A^{+\times}_{i,j}(\phi,\omega)
  + A^{\times+}_{i,j}(\phi,\omega)
\Big).
\end{align}

Where the bracket is the averaging of the $A^{++}$, the centered difference of both X and y. $A^{+ \times}$, the diangonal interaction. And $A^{\times +}$, another diangonal rotated relative to the other.

Using these principles and a lot of mathematics we gain these massive equations.
\begin{verbatim}
0 &=
\frac{\omega^{n+1}_{i,j} - \omega^{n}_{i,j}}{h_t}
+ \frac{1}{4}\left[
    A_{i,j}(\phi^{n+1}, \omega^{n+1})
  + A_{i,j}(\phi^{n},   \omega^{n+1})
  + A_{i,j}(\phi^{n+1}, \omega^{n})
  + A_{i,j}(\phi^{n},   \omega^{n})
\right] \nonumber \

\[1em]
&\quad + \frac{1}{4}\left[
    A_{i,j}(j^{n+1}, \psi^{n+1})
  + A_{i,j}(j^{n},   \psi^{n+1})
  + A_{i,j}(j^{n+1}, \psi^{n})
  + A_{i,j}(j^{n},   \psi^{n})
\right], \tag{86a}
\

\[1em]
0 &=
\frac{\psi^{n+1}_{i,j} - \psi^{n}_{i,j}}{h_t}
+ \frac{1}{4}\left[
    A_{i,j}(\phi^{n+1}, \psi^{n+1})
  + A_{i,j}(\phi^{n},   \psi^{n+1})
  + A_{i,j}(\phi^{n+1}, \psi^{n})
  + A_{i,j}(\phi^{n},   \psi^{n})
\right], \tag{86b}
\

\[1em]
\omega^{n+1}_{i,j} &=
- \Delta_x \phi^{n+1}_{i,j}
- \Delta_y \phi^{n+1}_{i,j}, \tag{86c}
\

\[1em]
j^{n+1}_{i,j} &=
- \Delta_x \psi^{n+1}_{i,j}
- \Delta_y \psi^{n+1}_{i,j}. \tag{86d}
\end{verbatim}


If you wish to see the derivation out, I would reccomend reading the paper.

\subsection{Discrete Conservation Laws}


Using the same methodology we can find our earlier conservation laws

Total energy
\[
E
= \frac{hx\,hy}{2}
  \sum_{i,j} \left( \phi_{i,j}\,\omega_{i,j} + \psi_{i,j}\,j_{i,j} \right)
= \text{const.}
  \]

  Magnetic Helicity
\[
C_{\mathrm{MH}}
= hx\,hy \sum_{i,j} \psi_{i,j}
= \text{const.}
  \]

  $L^2$ norm of $\psi$
\[
C_{L^2}
= hx\,hy \sum_{i,j} \psi_{i,j}^{\,2}
= \text{const.}
  \]

  And cross Helicity
\[
  C_{\mathrm{CH}}
= hx\,hy \sum_{i,j} \omega_{i,j}\,\psi_{i,j}
= \text{const.}
\]

\subsection{Numerical Methods}

So now we have out implicit forms tangled in non-linearity.

To solve this we use Newton's Method linearize the residual, solve the resulting linear system, update your guess, repeat until the residual is essentially zero. That linear system at each Newton step gets handed off to GMRES, an iterative Krylov-subspace solver, and to make GMRES converge fast, they build a custom, physics-informed preconditioner — following earlier work by Chacón and collaborators — that exploits the specific structure of the vorticity and Ohm's law equations.

The elliptic bits — inverting the Laplacian to get the varibles — get handled by whatever general-purpose elliptic solver you have on hand (they use PETSc); if your domain happens to be doubly periodic, there's an even cleaner trick available, since periodic boundary conditions mean the discrete Laplacian is exactly diagonalized by a Fast Fourier Transform, so you can invert it directly and exactly in Fourier space, no iteration required.


I mimic this in python by setting up a file to create the initial condition, the operators of the spectral solver and Arakawa's method, and integrator file to run it through the actual equations derived earlier.

\subsection{Validation}%
                       %
Here I make sure my model matches the paper through Orszag-Tang Vortex

(image)

Current sheet

(image)

Current sheet with Electron Inertia

(image)

Here the models match that of the paper, and matches within machine persicion the analytical solution.

\subsection{Random Test}

Now that we have a working model, let's use it for something cool. Here is Kinetic-Alfven Trublence

(movie here)



\section{Why: Developing Physics}


\subsection{Planing}

Essentially, one person teaching another parabolic tragectories but the student keeps asking why in a generally eurdite way. Analogous to Socratic Dialog.
\begin{enumerate}
  \item Introduction to Parabolic tragectories
\end{enumerate}


\subsection{Introduction}

Here we set the stage, a stage. An introductory physics teacher is having a one-on-one lesson with a young student. A highly intelligent student that wants to get to the base of everything. The most fundemental of the fundemental.


\section{Can you Destroy a Universe with a Force?}

\subsection{Introduction}

I am sure everyone knows of the extreme forces created by fictional characters. Blowing up planets, stars, galaxies, and even Universes. Now these obviously aren't physically realistic, but it begs the question, how much? What would such extreme forces actually look like?

\subsection{Planck Force}

To start with this, we will start with the theoretical highest force that still works with modern physics. The Planck Force. To explain this, the Planck Units are essentially normalizing and using units and physics starting with only the speed of light, gravitational constant, and Planck's Constant. Beyond the simple exersice of mathematical reasoning and the possibility of simpler mathematics. As is the case with other Natural Units, these units have been found to also coincide with when both quantum effects and gravity become significant. Making traditional theories fall apart, as no one knows how these two things will interact.

Using this method, the force unit becomes
\[
\frac{c^4}{G}
  \]
  Obviously an extremely large number, in newtons it is roughly two times ten to the power of 44.

  Many have conjectured that this is the largest possible force that can occur. Others have disputed. Others have found that according to traditional General Relativity it is simply a fourth of the Planck Force.

  Until a theory of Quantum Gravity has been resolved, no one will know.

\subsection{Striking the Earth with a Planck Force}

While obviously, without a theory of quantum gravity, no one really know what will happen with a Planck Force, but using General Relativity to guess should be fun.

I used Einstein's toolkit, a very helpful open source numerical general relativistic code base, and added a thorn to strike an object at a point with a force. Then I edited their Neutron Star simulation to be Earth adjacent, obviously not exact but with such extreme cases I doubt it would make much of a difference.

Doing this, was surprisingly difficult. I thought it would be the easiest part of this video, but I was wrong.

after running it I go this

(Insert movie)

While you likely can't read the numbers, and they are very strange units at that. So I will explain, very quickly the of contact very quickly explodes into a vacum at relativistic speeds. Then, as the relativistic mass adds up, gravity slowly overcomes the pressure wave. Collapsing it into a small point. A point roughly the size of a single mesh in the simulation. Thus, including things such as the likely fusing of particles, neutron degenerate matter, or if it would collapse into a black hole would require a more advanced simulation.

There is another problem. The force being so great would lead to particles reaching super-luminal speeds before the next time step. Leading to the simulation coming to errors as the mass become imaginary. The creators of the toolkit obviously came to a solution, but I am not exactly sure what their solution is. Thus, how accurate this simulation is, is knowledge beyond mine.

I tried to decrease the computational complexity and decrease the time-step substantially. I made the simulation symmetric so that I only had to simulate one eighth of the Earth, 2 dimensional rather than 3d, and increase the size of the mesh.

I eventually got to a single time step being one ten trillionth of a second. Still getting the same warning. So, I eventually had to accept the fact that I don't know. I would think that a collapse would happen but not into a black hole. Likely a mini-star that would eventually explode out, but that is something that I can't confirm.

To do so would require much more than my computer and messing with Einstien toolkit.

\subsection{Assumptions about a Universe}

While that was rather disappointing, there is still the question of the universe. Which is surprisingly much simpler.

Anyways, the we need to make assumptions about said universe.

One, our current observable universe, 93 billion light years in diameter, is the size we will use. Assuming that matter outside of such volume would have a negligible effect on the simulation.

Two, we assume that spacetime is mostly flat. Most observations put that the universe has an extremely small overall curvature, but the small change likely wouldn't change the results of the simulation much.

Three, take dark energy to be constant, such as a vacum energy model. While recent studies despute this, they have trouble coming up with exact values so for now I will leave it as is to simplify.

Four, with dark matter, we take the estimated values of both Hubble's estimations of dark matter and Shoe's as well.

Five, we take a relatively isotropic density in the universe. Meaning a relatively constant one, rather than variable.

Six, when we later talk about the energy of the force, assume that it is some kind of magic or from another universe. As it is ``creating'' or adding energy to the universe that was not there previously.

Now that we have these.

\subsection{Modeling a Universe}

Now, how to actually destroy a universe. Being filled with so much empty space it is pretty obvious you can't just punch it and blow it up. The small quarks in the plasma might hit a couple of things, but would be pretty useless and definitely wouldn't make it to other galaxies. So, that leads us to our solution. In the previous simulation, the Earth collapsed from gravity, so it is conceivable that a more intense force would collapse into a black hole.

So, first we work backwards. To get force we can take energy, and to get energy we can find the analogous mass needed to overcome dark energy.

Starting with the FRW Friedmann equation:
\[
H(a)^2= H_0^2 \left( \Omega_{r0}\,a^{-4} + \Omega_{m0}\,a^{-3} + \Omega_{\Lambda 0} \right)
\]
Here we will compress the parentheses into $E(a)$ for simplification.

Next, we take the Schwarzschild-de Sitter metric function for a relatively flat expanding space time.

\[
f(r) = 1 - \frac{2GM}{c^2r}-\frac{\Lambda r^2}{3}
  \]
  Here, one can take the derivative and find when that derivative is equal to 0 to find the point in which the universe becomes static to be

\[
    r_c = \left(frac{3GM}{\Lambda c^2}   \right)^{\frac{1}{3}}
  \]

  Then solving for mass when the radius is the current size of the universe we get
\[
M_{crit} = \frac{\Lambda c^2}{3G}R^3_{obs} = \frac{\Omega_{\Lambda 0}H^2_0}{G}R^3_{obs}
  \]

  Now for the exact solution we take

\[
ds^2 = -c^2dt^2 + \frac{R'(r,t)^2}{1+2E(r)/c^2}dr^2 + R(r,t)^2 d \Omega^2
  \]

  Here we take the Lemaitre-Tolman-Bondia metric, an exact solution for Einstein's equation for a lump of matter with no assumed symmetry.

  (show algebraic manipulation)

  Here we see that the structure of the shell makes no difference, only the total mass of it.

  From here simple algebra shows allows us to to find the energy needed to lead to an eventual collapse and the function that can be solved to find when it crosses.

  To actually solve, we need values. For the hubble constant and the mass and energy of the universe. Using the two values from Hubble telescope and Shoe, we use Bayes uncertainty to create a Monte Carlo simulation to find the list of values.

\subsection{Results}

From here, we run the simulation. Getting the below results

(results)

with the total energy needed to over come

\subsection{Interprating the Results}

Now, the purpose of this, wasn't to find the energy, but rather the force. I know most people watching this know the general work formula, but in the case of relativity, this doesn't always work. Though, with a couple of simplifications, we will end up with a just as simple formula.

We start with a practically flat spacetime to make the calculations simple.

From here, the fact of the conservation of energy and momentum, makes it so that the stress energy tensor, which has both energy and momentum is clear. Derived to be
\[
\nabla_\mu T^{\mu \nu}= f^\nu
  \]

  Where we take the contravairiant derivative of the stress energy tensor to get the force.

  where we can expand to
\[
f^\nu = \pdv{T^{0\nu}}{x^0} + \pdv{T^{1\nu}}{x^1} + \pdv{T^{2\nu}}{x^2} + \pdv{T^{3\nu}}{x^3}
  \]

  where we can take the first term to find that this is simply the energy added by the force by volume. Then the other three terms simplify to be

\[
f^i = \frac{1}{c}\pdv{T^{0i}}{t} + \partial_j T^{ij}
  \]

  which s essentially the change in momentum density plus the momentum flowing through boundary.

  To simplify the calculations, we will focus just on the energy density portion. The result would be the same just differing math. Say the change in energy is 3.6 times ten to the power of 71. This force is spread through a volume of 1 meter, energy density and thus force would be 3.6 times 10 to the power of 71 Newtons.

  As I have mentioned, this is an extreme simplification, but when working with energies such as these, there is really no physical object that would lead to the initial energy or the pressure gradient to have any non-negligible effect.

  So now we have the force. 36 followed by 70 zeros Newtons.

  About 67 orders of magnitude beyond the average punch of 1,000 Newtons. About 66 orders of magnitude above the highest ever recorded punch of ~69,000 Newtons.

  For comparasion, this is 25 orders of magnitude above the estimated energy release of the big bang.

  On the order of comics, the best feat that I could come up with a real number too, not estimations based on the wacky comic physics would be superman pulling a chain of planets as very quick speeds. With other people's estimates coming at $10^{38}$ to $10^{40}$ newtons. Still on the upper end being 31 orders of magnitude away. Humans are closer to the superman's feat than superman is to the calculated feat.


\subsection{General Questions}

There are a couple of questions that I can anticipate. If gravity moves at the speed of light, how would that changes things. The answer is not an insane amount. Red shifting and losing energy is only for waves, the constant gravity would not lose energy. Our model make sure the gravity overcomes dark energy by definition. The main change would be the change in the radius of the universe and the possibility of parts of the currently observable universe going past the cosmic horizen where gravity would never be able to reach it. The two aspects conteract one another, and since we are working with such large numbers I assume they wouldn't change much. Though, as always, I could always be wrong.


  Thank you for watching


\section{Over Engineering Celestial Mechanics: WHCKL model}

\subsection{Planning}%
\begin{enumerate}
  \item Introduction
  \item Scope
  \item Deriving WH
  \item WH Maping
  \item Deriving Symplectic correctors
  \item using variations
  \item Splitting method
  \item Lazy Kernel
  \item General Improvements
  \item Comparing to other methods
  \item video of solar system
\end{enumerate}

\subsection{Introduction}

(general animate introduction)

Hi everyone

\subsection{Scope}

Essentially the scope and goals of this project are to develop a celestial mechanics numerical model that is mathematically complex enough to be a challenge. To slightly constrain the scope, the goal will be to allow for highly efficient long term keplarian celestial mechanics. On the order of millions of even billions of years with relative accuracy.

After some research, the clear, for very long term celestial mechanics the clear modern winner is Wisdom-Holman

Then, when researching this program I found that REBOUND, an open source code for celestial mechanics found ways to improve its accuracy that introduced new fun challenges. Such as symplectic correctors, variational methods, splitting methods, and even something they call a lazy kernel. So, let us get into it.


\subsection{Deriving WH}

To derive Wisdom-Holman, we start with a traditional N-body hamiltonian in inertial cartesian coordinates

\[
H(q,p) = \sum_{i=0}^{N-1}\frac{P^2_i}
{2m_i} - G \sum_{i<j} \frac{m_i m_j}{\norm{q_i - q_j}} \]

With i = 0 to be the central star, and the rest being the planets.

From here, the hamiltonian is not sperable, so we make a canonical transformation into heilocentric canonical coordinates.

Where we define $Q_i = q_i -q_0$ where essentially Q is the position of the planets in refrence to the sun. The corresponding canonical momenta doesn't change. $P_i=p_i$.

With the exeption of the sun which changes by the transformation. $P_0 = p_0 + \sum^{N-1}_{i=0} P_i$.

Then we define capital M to be the cumulative mass

From here it is simply derived that the Hamiltonian becomes
\[
H = \frac{P_0^2}{2M} + \sum_{i=1}^{N-1} \left[\frac{P^2_i}{2m_i} - \frac{Gm_0 M_i}{Q_i}  \right] + \sum^{N-1}_{i=1} \frac{P^2_i}{2m_0} - G \sum_{i<j, i,j > 0} \frac{m_i m_j}{\norm{Q_i-Q_j}}
  \]

  Here we can divide the hamiltonian into $H = H_{sun} + H_{kepler} + H_{int}$.

  This is the traditional method; though, for later reasons we will get into, rather than Heilocentric coordinates we will use Jacobi, meaning in refrence to the center of mass rather than the central body.

  Where the center of mass of all particles interior to the i-th particle is $R_{i-1}$, thus their coordinates become $r'_i = r_i - R_{i-1}$.

  Because jacobi coordinates are linear functions of cartesian, the rest of the quantities, such as velocity and acceleration change in the same way. Momenta is different however.
\[
p'_i = m'_i \dot{r'}
  \]
  where the reduced mass is equavlent to.
\[
m'_i = m_i \frac{M_i-1}{M_i}
  \]



  Using these coordinates the N-body hamiltonian instead becomes
\[
H = \sum_{i=0}^{N-1}\frac{p'^2_i}{2m_i'} - \sum_{i=0}^{N-1} \sum_{j=i+1}^{N-1} \frac{G m_i m_j}{\norm{r_i - r_j}}
  \]

  Keeping the potential term in cartesian coordinates and only writing the kinetic in jacobi.

  Next, for the sake of the algebra we add and subtract the term
\[
H_{\pm} = \sum_{i=1}^{N-1} \frac{G m'_i M_i}{\norm{r'_i}}
  \]

  This grouping leads to the Hamiltonian





\[
H =
\underbrace{\frac{(p_0')^{2}}{2 m_0'}}_{H_0}
+
\underbrace{
\sum_{i=1}^{N-1}
\left(
\frac{(p_i')^{2}}{2 m_i'}
-
\frac{G\, m_i' M_i}{\lvert r_i' \rvert}
\right)
}_{H_{\text{Kepler}}}
+
\underbrace{
\sum_{i=1}^{N-1}
\frac{G\, m_i' M_i}{\lvert r_i' \rvert}
-
\sum_{i=0}^{N-1}
\sum_{j=i+1}^{N-1}
\frac{G\, m_i m_j}{\lvert r_i - r_j \rvert}
}_{H_{\text{Interaction}}}.
\]







Here one can further split the kepler term into

\[
    (H_{ekpler})_i = \frac{p^{'2}_i}{2m'_i} - \frac{Gm'_i M_i}{\norm{r'_i}}
  \]
  each Hamiltonian describes the keplarian motion of the i-th particle around the center of mass of all interior  particles.

  After some more algebra we can split the jacobian into two parts, one that can be easily computed in cartesian and the other in Jacobi coordinates

\[
H_{int} == \sum_{i=2}^{N-1} \frac{Gm'_i M_i}{\norm{r'_i}} - \sum_{i=0}^{N-1} \sum_{j=i+1}^{N-1} \frac{Gm_im_j}{\norm{r_i - r_j}}
  \]

  We now have three speratable equations that a all much easier to integrate than the highly non-linear one we stated with. Where $H_0$ is motion along a straight line, $H_{Kepler}$ is the hamiltonian for kepler orbits, and $H_{int}$ is the interaction of particles outside of the i-th particle. To note, this methodology only works for near keplarian orbits, in say a three body binary this method would be quite unphysical quite quickly.

\subsection{WH Maping}

We do this splitting as the origional equation's equation of motion have no analytical solutions.

\[
H(q,p) = \sum_{i=0}^{N-1}\frac{P^2_i}
{2m_i} - G \sum_{i<j} \frac{m_i m_j}{\norm{q_i - q_j}}
  \]

  Thus, approximate solutions are needed. Many methods exist but the method we have been developing is obviously splitting the Hamiltonian into smaller parts that can be integrated easily.

  Analytical solutions can be found for the evolution of the system under each individual Hamiltonian.


  Thus, to create the symplectic integrator using operator split method.

  To do this we start with the familiar Liouville operator
\[
L_H f = \{f, H\}
  \]
  If you are not familiar, this is simply the same as saying that
\[
\dot{f} = \{f,H \}
  \]

  This obviously looks like an ordinary differential equation that has the formal solution of

\[
f(t+h) = e^{h L_H}f(t)
  \]

  Because this equation is for a differential operator not a variable, this exponential opperator is instead defined through:
\[
e^{hL_h} = I + hL_H + \frac{H^2}{2!}L^2_H + ...
  \]
  where the key difference is that differential operators do not commute as traditional variables do. This here is the backbone of the symplectic integrator, to approximate this non-commutable equation.

  Using this we can take our origional Hamiltonian in Poisson form to be:

\[
L_H = L_0 + L_{Kep} + L_{int}
\]

let us turn those into A, B, and C:
\[
L_H = A + B + C
\]
Then add A and B to make a and turn C into B
\[
L_H = A + B
\]
then the exact propegator is obviously
\[
e^{h(A+B)}
\]
Now that we have that we can use the generic strang method.

\[
\boxed{
e^{\Delta t(A+B)}
\approx
e^{\frac{\Delta t}{2}A}
e^{\Delta t B}
e^{\frac{\Delta t}{2}A}
}
\]

Or the drift-kick-drift. Drifting along the kepler motion and motion of the center of mass and kicking using the inteaction between planets.

\subsection{Symplectic Correctors}

Another method we will add are symplectic correctors. A symplectic corrector is simply a small adjustment to the coordinates that one does before and after running a symplecitc integrator so that the numerical trajectory more closely matches the real physical trajectory.

Because they are only applied at the very beginning and end, one can add arbitrarily higher order correctors with negligible effects on the speed.

We choose an eleventh order corrector that effectively brings the effective order of the integrator from second order to fifth, while barely changing the time it takes to run the program.

The exact benefits of the correctors change. As in Earth mass planets, beyond a fifth order corrector there is little change, but in Jupiter mass ones an eleventh order method is quite helpful.


\subsection{Building the Code}

When building this code, taking optimization to heart I used FORTRAN. A computational science coding language that is very fast. Optimized it is comparable and arguably faster than optimized C++, it also has the added benefit of being much easier to read.

I found a paper that already derived the best ways to optimize the subcalculations, specifically coordinate transformations between Jacobi and Cartesian, Newtons Method, c-functions, and a couple of other things.

I won't bore you with the details, but I eventually added the fiften of the most massive bodies that orbit the sun. The eight planets, Ceres, Pluto, Haumea, Makemake, Gonggong, and Erirs. Yes, those are real names of dwarf planets.

I didn't add moons as this simulation, by using the kepler hamiltonian, assumes that the body orbits the center of mass of interior object, but moons obviously don't do that. Such, the simulation is not as physical for say the Earth whose moon is a significant compared to its mass. But it still works as a fun project with beyond that being very accurate.

After that I set the time step to roughly 18 days, the max speed while still being accurate. I did underestimate how many years the simulation was able to complete so Mercury did become unbounded near the end

\subsection{Results}

Here we now come to the fruits of our labor. I really should have worked on another program to have something to compare to, as these numbers are pretty meaningless on their own. Just that they are about as efficient as it gets...

Well, the purpose was the process not the result

\subsection{Watching it}
Though, to get a little extra out of it, let us watch the animation of a couple of the planets.


\section{Lie Operator Splitting: One of the ten Pillars}

\subsection{Introduction}

Hi, welcome to a semi-advanced Introduction into Lie Operator splitting, a mathematical Approach.


\subsection{Scope}

Essentially, the scope of this video being an introduction into Lie Operator splitting, a method that I have used several times in past videos.

A good introduction to its importance is a quote by Dr. Strang, ``If, as has sometimes been argued, there are only ten big ideas in numerical analysis and all the rest are merely variations on those themes, splitting is undoubtedly one of them.'' as it is at the foundation of so much of numerical analysis. I primarily have used it for symplecic integrators, but its uses are far beyond that. Other structure preserving methods, statistical mechanics, nearly integratable systems, optimization, quantum mechanics, and really anything that can be described by differential equations without easily found analytical solutions.

n fact, one could say that, at least since Descartes stated his four rules of logic in the Discours de la méthode, the notion of sub-dividing a complicated problem into its simpler constituent parts, solving each one of them separately and combining those separated solutions in a controlled way to get a solution to the original overall problem constitutes a guiding principle in all areas of science and philosophy.

\subsection{Motivating Operator Splitting}

As in most cases in applied mathematics, the easiest start to defining new mathematics is to motivate their existence. To confine this agumentation, let us look at the simpliest case of a non-integratable Hamiltonian. An Anharmonic oscillator defined through this equation:
\[
H(q,p) = \frac{p^2}{2m} + \frac{1}{2}m \omega^2q^2 + \lambda q^4
  \]

  Here, this equation's derivatives have no elementary closed-form solution

  But we can split it into its potential and kinetic energy
\[
H_T = \frac{p^2}{2m}
  \]
\[
H_V = \frac{1}{2}m \omega^2q^2+ \lambda q^4
  \]

  Here these problems have trivial solutions.

  Thus, if one wishes to solve them together one can trivially see that the Hamiltonian vector field is
\[
X_H = X_T + X_V
\]
(show vector fields adding together)

The corresponding Liouville operators are
\[
L_H = L_T + L_V
\]

Then if we introduce a phase-space vector

\[
    z =
  \begin{pmatrix}
    q \\
    p
  \end{pmatrix}
\]

Solving the differential equation we get
\[
z = e^{h(L_T + L_V)}
  \]

  Solving this may seem easy. Your first instinct would probably be solve the inside the parenthasies, but this is obviously the same thing as integrating the whole, making the earlier splitting useless.

  Another methodology would be through commuting
\[
e^{hL_T}e^{hL_V}
\]

But, as this doesn't work as they do not commute.
\[
[L_T,L_V] \neq 0
  \]

  We can see this directly for solving for q
\[
[L_T, L_V]q = -\frac{V'(q)}{m} - 0
  \]

  Thus this take naive one-step composition to be problematic.

  I will note, this traditional splitting can work, it is only first order accurate and has a local error of $O(h^2)$

\subsection{Deriving Strang Splitting}

Taking our differential equation earlier, replacing T and V with A and B:

\[
\dot{z} = (A + B)z
  \]

  Using taylor expansion we can expand to:
\[
z(t+h) = \left[I + h(A+B) + \frac{h^2}{2}(A^2 + AB + BA + B^2) \right]z + O(h^3)
  \]

  Thus, any second order method must reproduce the middle section.

  To do that, let us assume a three state composition that has factors a and b
\[
S(h)= e^{ahA}e^{bhB}e^{ahA}
\]

We can thus expand the exponentials to
\[
e^{ahA} = I + ahA + \frac{a^2h^2}{2}A^2+O(h^3)
  \]
\[
    e^{bhB} = I bhB + \frac{b^2h^2}{2}B^2+ O(h^3)
  \]

  Multipliying together we get
\[
S(h) = I + h(2aA +bB) + O(h^3)
\]

Where as we previously noted that the exact solution should be
\[
I + h(A+B)+ O(h^3)
\]

Therefore, a should equal one half while b should equal 1.

Creating a second order method with a local error of $O(h^3)$.

An importaint thing to note, while if you assume that the lower case a's could have differing values, we gain infinitely many solutions, but these solutions lose symmetry that would then lead to another first order equation.

\subsection{Adjoint formulism}
self adjoint formulisms are very simple, they are when a numerical method satisfies
\[
\psi_h = \psi^{\star}_h = (\psi_{-h})^{-1}
  \]

  Where psi star is the inverse composition of psi. ie if $\psi = x^1 \circ x^2$ then $\psi^{\star} = x^2 \circ x^1$.

These self-adjoin is sometimes called time-symmetric as ...
\subsection{Higher Order Splitting and Composition}
``it is a meta-theorem of numerical analysis that second order methods often achieve the right balance between accuracy and complexity.'' - Strang

While most often 2nd order is used, it isn't the only one. There are methods for developing higher order methods, particularly used many places. The most obvious being simulations of the solar system.

So, if we start with the strang splitting
\[
S^{[2]}_h = \varphi_{h/2} \circ \varphi_{h} \circ \varphi_{h/2}
  \]

  The composition is then

\[
\psi_h = S_{\gamma_s h} \circ S_{\gamma_{s-1} h} \circ ... \circ S_{\gamma_1 h}
  \]

  The order is atleast three if
\[
\sum^s_{j=1} \gamma_j = 1
  \]
\[
\sum^s_{j=1} \gamma_j^3 = 0
  \]

  The smallest number of j that satisfies this equation is three. By imposing a symmetry where $\gamma_3 = \gamma_1$, one can relatively easily derive a solution for $\gamma_1$ and $\gamma_2$.

Futher, in general one can take the equation
%\[
%S_{h}\!\left[2^{k}\right]= S_{\gamma_{1}\, h\!\left[2^{k-2}\right]}\circ S_{\gamma_{1}\, h\!\left[2^{k-2}\right]}\circ S_{\left(1 - 4\gamma_{1}\right)\, h\!\left[2^{k-2}\right]}\circ S_{\gamma_{1}\,h\!\left[2^{k-2}\right]}\circ S_{\gamma_{1}\, h\!\left[2^{k-2}\right]}
%  \]

  To find arbitrarily high orders of order 2k where $k \geq 2$.

There are many other ways to find higher order splitting, a more detailed observation of these will be in the description.



\section{Discontinuous Galerkin Method: Self Gravitating Euler Equations}
\subsection{Planning}
\begin{itemize}
  \item Introduction
  \item Scope
  \item Motivating DG
  \item General Introduction to DG
  \item Physical Model Model
  \item Deriving Projection
  \item Structure preserving LDG
  \item Time Discretization
  \item Ossilation Elimination
  \item Structure Preserving Properties
  \item The Code itself(Flowchart)
  \item Results
\end{itemize}
\subsection{Introduction}

Hi, welcome to Discontinuous Galerkin Method: a focus on Self-Gravitating Euler Equation


\subsection{Scope}

Here we address the scope of this video. It is meant as an introduction into DG, using an example to help facilitate understanding. For both the purposes of challenging myself, I have chosen a particularly challenging one I have recently read about.

Taking a self-gravitating fluid and using the structure preserving capacity of DG to preserve positivity, ie that pressure, mass, energy, etc are all positive at all points, and well balanced, ie compressible fluid that admits hydrostatic equilibrium. For this video I won't go to far into the math of these more advanced subjects.

Also, for the rest of the video I will refer to the DG are ``DG''

\subsection{Introduction to DG}

Before we get into the math, we will take a more intuitive look before we get into the math. A good way to start is to look at the whole family tree. I assume anyone who is watching this has heard of and understands finite difference models; approximate spacial derivatives with stencils on point values. Further, we get to finit volume which trackes cell averages and evolve them via fluxes computed at cell faces from a Riemann solver. Then we come to continuous finite element or Galerkin who represents the solution as one global piecewise-polynomial function, continuous across element boundaries. Higher order and structural preserving methods are baked into the equations. The separation between DG and CG is that these polynomials are discontinuous, exactly as you would expect. The reason for this is that in many cases, the physical situations are partially discontinuous. Think of an explosion in a fluid, it takes time for the explosion to effect outer areas of the fluid, more time than would be found through a continuous polynomial.

essentially, we divide the space into elements
\[
\Omega = \cup_e \Omega_e
  \]

  We approximate the solution to the differential equation using polynomials
\[
u_h(x,t) \sum_{j=0}^N u_j(t) \phi_j(x)
  \]

  Next we write the PDE in its weak form
\[
  \pdv{u}{t} + \grad \cdot \mathbf{F} = S(u)
\]


Multiply the test function $\phi$ then integrate by parts

\[
\int_\Omega \phi \pdv{u}{t} dV - \int_\Omega \grad \phi \cdot \mathbf{F} dV + \int_{\partial \Omega} \phi \mathbf{F} \cdot \mathbf{n} ds = S(u)
  \]

  They communicate between elements via fluxes, that can be defined via your physical system

  Then we can generalize the fluxes into the residuals from the weak form equation used previously used to rearange into

  \begin{equation}
\mathbf{R}(\mathbf{U})
=
\mathbf{R}_{\mathrm{vol}}
+
\mathbf{R}_{\mathrm{surf}}
+
\mathbf{R}_{\mathrm{src}}
\end{equation}

\begin{equation}
\mathbf{R}_{\mathrm{vol}}
=
\int_K
\nabla\boldsymbol{\phi}
\cdot
\mathbf{F}(\mathbf{u}_h)
\,dV
\end{equation}

\begin{equation}
\mathbf{R}_{\mathrm{surf}}
=
-\int_{\partial K}
\boldsymbol{\phi}
\,\mathbf{F}^{*}
\,dS
\end{equation}

\begin{equation}
\mathbf{R}_{\mathrm{src}}
=
\int_K
\boldsymbol{\phi}
\,\mathbf{S}(\mathbf{u}_h)
\,dV
\end{equation}

\begin{equation}
\mathbf{R}(\mathbf{U})
=
\int_K
\nabla\boldsymbol{\phi}
\cdot
\mathbf{F}(\mathbf{u}_h)
\,dV
-
\int_{\partial K}
\boldsymbol{\phi}
\,\mathbf{F}^{*}
\,dS
+
\int_K
\boldsymbol{\phi}
\,\mathbf{S}(\mathbf{u}_h)
\,dV
\end{equation}

\begin{equation}
M\frac{d\mathbf{U}}{dt}
=
\mathbf{R}(\mathbf{U})
\end{equation}

from here, using the basis functions from the mesh we get the mass matrix which is the matrix representation of the inner product on your finite-dimensional approximation space.

From here, the derivative of the functions as a function of time can be derived via
\[
\pdv{u}{t} = M^{-1}rhs(u_h)
  \]

  Now, listening to all of this may seem fairly confusing, this is doubled by my admitted brevity, but I allow this brevity as I am going to take a concrete example as mentioned previously.

\subsection{Physical Model}

As mentioned earlier, we take the Euler-Poisson for a self-gravitating fluid, such as a toy model of a star. While keeping the simplicity of the model, I add a personal challenge by adding constraints of energy conservation and positivity. Using the general Euler-Poisson equation
\[
\text{insert long equation here}
  \]

  Then adding Polytropic steady-states and Lane-Emden equations in order to enforce EOS and positivity.

\[
\text{equation}
  \]

  I know I once again was very brief, but that is simply because the actual equations don't matter too much, just the general shape and view of them so that we can create the code.

\subsection{Mesh and Projection}


We start by assuming a uniformly space square/cube meshing

(image of square meshing)

We can then define $n^+, n^-$ as the normal of the flux.

Next to define the projection, but to do that we define our test problem. Simply a weighting function within the polynomial space that allows us to integrate more easily.

(show examples in traditional methods)

Then we can define the projection matrix P via the $L^2$ projection:
\[
\int_K u v dx = \int_K \mathbf{P}(u) v dx
  \]

  Now we can start moving backwards by rearanging the equation to get
\[
P(u)(x) = \sum c_j v_j(x)
  \]

  By defining the mass matrix as
\[
M_{ij} = \int_k \phi_i v_j
  \]

  where phi of i is the basis function, we use the nodal basis function, coefeciencts fairly easily looked up.

  we can also define the right side vector as
\[
b_j = \int_K u v_j dx
  \]

  We can rearagne our origional projection function
\[
P(u)(x) = \sum (M^{-1}b)_j v_j(x)
  \]

\subsection{Flux}

We start with a basic peicewise function for flux of a fluid.
\[
    F^{hllc}(U_L, U_R; \mathbf{n}) =
\begin{cases}
  \mathbf{F}(U_L), \text{if } 0 \leq S_L, \\
  \mathbf{F}_{\star L}, \text{if } S_L \leq 0 \leq S_\star, \\
  \mathbf{F}_{\star R}, \text{if } S_\star \leq 0 \leq S_R, \\
  \mathbf{F}(U_R), \text{if } 0 \geq S_R, \\
\end{cases}
\]

This is the standard HLLC or Harten–Lax–van Leer–Contact flux solver. As this is more an introduction into DG than fluid dynamics I won't go through and expand all of these out, but essentially, it approximates the Remanian problem to approximate the flux of conserved quantities like energy, momentum, mass, etc.

I also slightly modify it in order to enforce Well-balencedness by adding the earlier Lane-Emden equation to the flux equation.

\subsection{Scheme}
Now building the actual numerical method. Before we get to the well-balanced, energy-conserving version, I want to build the plain scheme first, so that when I show you exactly how it fails near equilibrium later, you can point to the specific term responsible, rather than just taking my word for it.

Let's write out our conserved variables, flux, and source term the way we did generically before, but now specifically for the self-gravitating Euler system in $d$ spatial dimensions:
\[
\mathbf{U} =
\begin{bmatrix}
\rho \\ \rho \mathbf{u} \\ E
\end{bmatrix}, \qquad
\mathbf{F}(\mathbf{U}) =
\begin{bmatrix}
\rho \mathbf{u} \\ \rho \mathbf{u}\otimes\mathbf{u} + p\mathbf{I}_d \\ (E+p)\mathbf{u}
\end{bmatrix}, \qquad
\mathbf{S}(\mathbf{U}, \grad\phi) =
\begin{bmatrix}
0 \\ -\rho\grad\phi \\ -\rho\mathbf{u}\cdot\grad\phi
\end{bmatrix}.
\]

Here $E = \rho e + \tfrac12 \rho \abs{mathbf{u}}^2$ is the usual total (non-gravitational) energy, $e$ is the specific internal energy fixed by an equation of state, and for an ideal gas $p = (\gamma-1)\rho e$. With this notation, our full model compresses down to
\[
\mathbf{U}_t + \grad\cdot\mathbf{F}(\mathbf{U}) = \mathbf{S}(\mathbf{U},\grad\phi), \qquad \Delta\phi = 4\pi G \rho.
\]

$\phi$ is the gravitational potential and $G$ the gravitational constant. That second equation, Poisson's equation, is the entire reason this problem is harder than an ordinary Euler solve: it isn't fixed but rahter generated by the fluid, it changes shape every time the density field changes, and every single timestep we need to re-derive a self-consistent gravitational field from whatever the density happens to look like at that instant.

\subsubsection{Why local DG, specifically}

There are two broad strategies exist for handling the gravity term numerically. One is to treat the Poisson equation as a pseudo-hyperbolic problem, giving it an artificial wave speed and a relaxation time so the same Riemann-solver machinery built for the Euler equations can be reused directly on the gravity sector too. The other, which is the one this scheme takes, is to solve the elliptic Poisson equation directly. Direct solves used to be avoided in this setting because they historically meant either an FFT-based Poisson solver (fast, but restricted to very particular, usually periodic, domain shapes) or a multipole expansion (flexible, but expensive and awkward to make high order). Local DG sidesteps both complaints: it was originally invented specifically to let a DG method handle equations with second derivatives without ever forming a second derivative explicitly, and it inherits all of the DG machinery's flexibility with boundary conditions and geometry for free. That's the actual technical justification for reaching for LDG here rather than, say, a finite-difference Poisson solver bolted onto a DG hydro code.

\subsection{Standard LDG}

So concretely: how do we get $\phi$ out of a DG method when Poisson's equation has a second derivative in it, and DG's whole weak-form machinery is built assuming one derivative per equation? This is exactly the ``Local'' in Local DG. Following Cockburn and Shu's original construction, we introduce an auxiliary variable standing in for the gravitational force itself,
\[
\mathbf{g} = \grad \phi,
\]
and rewrite the whole system as a genuinely first-order system, no second derivatives anywhere:
\[
\begin{cases}
\mathbf{U}_t + \grad\cdot\mathbf{F}(\mathbf{U}) = \mathbf{S}(\mathbf{U}, \mathbf{g}), \\
\mathbf{g} = \grad\phi, \\
\grad\cdot\mathbf{g} = 4\pi G \rho.
\end{cases}
\]

Now every equation is first order, so we throw the exact same weak-form recipe from earlier at each piece independently. We work on a mesh $\mathcal{T}_h$ of elements $K$, with $\mathcal{E}_I$ the interior faces and $\mathcal{E}_B$ the domain boundary. On a shared face $\mathcal{E}$ between two elements $K^+$ and $K^-$, with outward normals $\mathbf{n}^+,\mathbf{n}^-$, define the usual average and jump of a scalar function $u$ and a vector function $\mathbf{q}$:
\[
\{\!\{u\}\!\} = \tfrac12(u^+ + u^-), \qquad [\![u]\!] = u^+\mathbf{n}^+ + u^-\mathbf{n}^-,
\]
\[
\{\!\{\mathbf{q}\}\!\} = \tfrac12(\mathbf{q}^+ + \mathbf{q}^-), \qquad [\![\mathbf{q}]\!] = \mathbf{q}^+\cdot\mathbf{n}^+ + \mathbf{q}^-\cdot\mathbf{n}^-.
\]

On a single element $K$, testing the three equations against $v \in V_h^k$, $\mathbf{w} \in \Sigma_h^k$, and $\psi \in V_h^k$ respectively (piecewise polynomial spaces of degree up to $k$, scalar and vector valued), and integrating by parts exactly the way we did for the single scalar conservation law earlier, we get:
\[
\int_K (\mathbf{U}_h)_t \, v \, d\mathbf{x} - \int_K \mathbf{F}(\mathbf{U}_h)\cdot\grad v \, d\mathbf{x} + \sum_{\mathcal{E}\in\partial K} \int_{\mathcal{E}} \widehat{\mathbf{F}}\, v^{\mathrm{int}(K)}\, ds = \int_K \mathbf{S}(\mathbf{U}_h,\mathbf{g}_h)\, v \, d\mathbf{x},
\]
\[
\int_K \mathbf{g}_h \cdot \mathbf{w}\, d\mathbf{x} = \sum_{\mathcal{E}\in\partial K}\int_{\mathcal{E}} \widehat{\phi}_h \, \mathbf{w}^{\mathrm{int}(K)}\cdot\mathbf{n}_{\mathcal{E},K} \, ds - \int_K \phi_h \grad\cdot\mathbf{w}\, d\mathbf{x},
\]
\[
\int_K \mathbf{g}_h \cdot \grad\psi\, d\mathbf{x} = \sum_{\mathcal{E}\in\partial K}\int_{\mathcal{E}} \psi^{\mathrm{int}(K)} \, \widehat{\mathbf{g}}_h\cdot\mathbf{n}_{\mathcal{E},K}\, ds - \int_K 4\pi G \rho_h \, \psi\, d\mathbf{x}.
\]

The superscript ``int$(K)$'' just means: take the value from inside the element $K$ we're currently integrating over, as opposed to the value carried over from the neighbor on the other side of the face (``ext$(K)$''). This distinction matters constantly in DG because our solution is discontinuous across faces by construction, so ``the value at the face'' is genuinely ambiguous until we pick a numerical flux to resolve it.

The first equation is exactly the fluid solve we already built, with $\widehat{\mathbf{F}}$ the HLLC flux from the previous section, evaluated using $\mathbf{U}_h^{\mathrm{int}(K)}$ and $\mathbf{U}_h^{\mathrm{ext}(K)}$ as the left and right Riemann states. The other two equations are the ``elliptic half'' of the scheme: a mixed formulation of Poisson's equation.

\subsubsection{Fluxes for the elliptic part}

Here's a subtlety that trips people up the first time they see LDG for an elliptic problem: gravity, unlike a fluid wave, doesn't propagate in a preferred direction. There's no upwind side to a Poisson equation the way there's an upwind side to a shock. So HLLC, or any Riemann-solver-based flux, simply doesn't apply here. Instead, following the standard LDG-for-elliptic-problems recipe, we use plain averages of $\phi_h$ and $\mathbf{g}_h$ across the face, corrected by a jump-penalty term that keeps the resulting linear system well posed:
\[
\widehat{\mathbf{g}}_h = \{\!\{\mathbf{g}_h\}\!\} - C_{11}[\![\phi_h]\!] - \mathbf{C}_{12}[\![\mathbf{g}_h]\!], \qquad
\widehat{\phi}_h = \{\!\{\phi_h\}\!\} + \mathbf{C}_{12}\cdot[\![\phi_h]\!],
\]
with $C_{11}$ a scalar penalty and $\mathbf{C}_{12}$ a vector penalty, both free parameters we get to choose. On a tensor-product (rectangular) mesh, a standard, well-behaved choice is $C_{11}=1$ with $\mathbf{C}_{12} = (0.5,0.5)^T$ in 2D or $(0.5,0.5,0.5)^T$ in 3D.

If you assemble the full linear system for $(\mathbf{g}_h, \phi_h)$ implied by these two weak forms, it has a block, saddle-point structure,
\[
\begin{bmatrix} K & B \\ -B^T & S \end{bmatrix}
\begin{bmatrix} \mathbf{g}_h \\ \phi_h \end{bmatrix}
=
\begin{bmatrix} g \\ f \end{bmatrix},
\]
which is annoying to solve directly because it's not symmetric. The standard move is a Schur complement: eliminate $\mathbf{g}_h$ algebraically to get a smaller system purely in $\phi_h$,
\[
A\phi_h = \left(S + B^T K^{-1} B\right)\phi_h = f + B^T K^{-1}\mathbf{g}_h,
\]
and now $A$ is symmetric positive definite, which means a Cholesky factorization or a (preconditioned) conjugate-gradient solve both work reliably. Since our mesh never moves or changes over the course of the simulation, that Cholesky factorization can be computed exactly once, at the very start of the run, and reused at every single timestep afterward — the sparsity pattern and the matrix entries here depend only on geometry and the penalty parameters, never on the current fluid state. (For a periodic domain the block structure changes slightly, the reduced system loses its symmetry, and an LU factorization is used instead — same idea, slightly different linear algebra.) This is a genuinely nice practical property of building gravity this way: the expensive linear-algebra setup work happens once, not once per timestep.

Initial conditions for the fluid variables are set by the $L^2$ projection $\mathbf{P}$ we derived earlier: $\rho_h^0 = \mathbf{P}\rho^0$, $(\rho\mathbf{u})_h^0 = \mathbf{P}(\rho\mathbf{u})^0$, $E_h^0 = \mathbf{P}E^0$. Wrapping the whole thing in a forward-Euler step in time gives a complete, working numerical scheme:
\[
\int_K \left(\frac{\mathbf{U}_h^{n+1}-\mathbf{U}_h^n}{\Delta t}\right) v\, d\mathbf{x} - \int_K \mathbf{F}(\mathbf{U}_h^n)\cdot\grad v\, d\mathbf{x} + \sum_{\mathcal{E}\in\partial K}\int_{\mathcal{E}} \widehat{\mathbf{F}}^n \, v^{\mathrm{int}(K)}\, ds = \int_K \mathbf{S}(\mathbf{U}_h^n,\mathbf{g}_h^n)\, v\, d\mathbf{x}.
\]

Density, momentum, and energy evolve; gravity is re-derived from the freshly updated density at the top of the next step, and the loop repeats. This is what I'll call the \emph{standard scheme} for the rest of the video. Let me be fair to it: this is a real, working, high order accurate method. Run it on a smooth manufactured solution and you'll see the convergence rates you'd expect from the polynomial degree you chose. Run it on a genuine shock and, with an appropriate limiter, it captures the shock. It is not a broken scheme.

It does, however, have one specific, well-defined blind spot, and that blind spot is the entire reason the rest of this video exists.

\subsection{Structure preserving DG}

A huge fraction of astrophysically interesting configurations spend almost all of their time sitting extremely close to hydrostatic equilibrium: gravity pulling material inward, pressure pushing back outward, velocity essentially zero, the two forces balanced to enormous precision. That exact balance is
\[
\mathbf{u} = \mathbf{0}, \qquad \grad p = -\rho \grad \phi, \qquad \Delta\phi = 4\pi G\rho.
\]

A common and analytically tractable family of such equilibria is the \emph{polytropic} equilibrium, where pressure and density are locked together by $p = \kappa \rho^\gamma$ for some positive constant $\kappa$. Substitute this and the hydrostatic balance into Poisson's equation and, after the classic Lane-Emden change of variables $\rho = \lambda\theta^n$ with $\gamma = (n+1)/n$, you land on the Lane-Emden equation. In spherical symmetry, with $r$ the (rescaled) radius, it reads
\[
\frac{1}{r^2}\frac{\partial}{\partial r}\left(r^2 \frac{\partial\theta}{\partial r}\right) = -\theta^n.
\]
This has closed-form solutions only for a few special integer values of $n$: $\theta_0(r) = 1 - r^2/6$ for $n=0$ (an incompressible, $\gamma=\infty$ fluid); $\theta_1(r) = \sin(r)/r$ for $n=1$ (a $\gamma=2$ fluid, and the one this video's numerical tests will actually use, since it's smooth, non-trivial, and exact); and $\theta_5(r) = 1/\sqrt{1+r^2/3}$ for $n=5$ ($\gamma=6/5$). For any of these, the full equilibrium profile is recovered from
\[
\rho^e = \lambda\left(\theta(\mathbf{x})\right)^n, \qquad \mathbf{u}^e = \mathbf{0}, \qquad p^e = \kappa(\rho^e)^\gamma, \qquad \phi^e = -\frac{\kappa\gamma}{\gamma-1}(\rho^e)^{\gamma-1}.
\]
This is the analytic ``star in a box'' that our scheme will be asked to hold perfectly still.

Here's the actual failure mode, stated precisely. The standard scheme discretizes the pressure-gradient contribution to the momentum flux and the gravity source term $-\rho\grad\phi$ using two completely independent numerical procedures: one is a DG surface/volume flux integral, the other is a DG volume source integral. There is no algebraic reason these two, separately truncated approximations should cancel exactly, even when the exact continuous quantities they're approximating do cancel exactly. What's left behind is $O(h^{k+1})$ discretization error sitting right on top of the equilibrium, at every single timestep. If the physical phenomenon you actually care about — a stellar oscillation mode, a slowly growing instability — is smaller in amplitude than that leftover truncation error, it's simply drowned out, and the only fix available to the standard scheme is to keep refining the mesh until the noise floor drops below your signal, which gets expensive fast, especially in 3D. This is the well-balanced (WB) problem, and a scheme that avoids it is called well-balanced.

There is a second, related failure worth separating out clearly, because it has a genuinely different cause: total energy conservation (TEC). Because gravity here is self-generated rather than externally imposed, the gravitational potential energy is honestly part of the system's total energy budget, and near equilibrium it's typically a large negative number nearly cancelling a large positive kinetic-plus-internal energy. Any small, systematic energy leak in the numerical scheme is no longer a rounding-error-level nuisance; it can be comparable to, or larger than, the actual dynamical energy of the perturbation you're trying to resolve, especially over long integrations or during genuine gravitational collapse, where nonlinear growth will happily amplify a small numerical energy leak into a qualitatively wrong answer.

We want both properties simultaneously, and it turns out they require two largely independent fixes.

\subsubsection{Fix one: splitting the potential for well-balancedness}

The core idea is to split both density and gravitational potential into a known, static equilibrium piece and an evolving perturbation piece:
\[
\rho = \rho^e + \rho^\delta, \qquad \phi = \phi^e + \phi^\delta, \qquad \mathbf{g}^e = \grad\phi^e,\quad \mathbf{g}^\delta = \grad\phi^\delta,
\]
with $\rho^e, \phi^e$ computed once at initialization from the analytic equilibrium above (via the $L^2$ projection: $\rho_h^e = \mathbf{P}\rho^e$, and $\phi_h^e$ from actually solving the discrete Poisson problem for $4\pi G\rho_h^e$) and then frozen, unchanged, for the entire simulation. Only the perturbation fields $\rho_h^\delta,\phi_h^\delta$ are recomputed every timestep. On its own this split buys us nothing; the trick is entirely in how the numerical flux and source term are built around it.

\paragraph{The modified HLLC flux.} Recall the contact property of HLLC: feed it two states at rest with equal pressure, $\mathbf{U}_L = (\rho_L, \mathbf{0}, p/(\gamma-1))^T$ and $\mathbf{U}_R = (\rho_R, \mathbf{0}, p/(\gamma-1))^T$, and it collapses exactly to
\[
\mathbf{F}^{hllc}(\mathbf{U}_L,\mathbf{U}_R;\mathbf{n}) = (0,\ p\,\mathbf{n}^T,\ 0)^T.
\]
You can check this directly from the definitions: with zero velocity on both sides, the estimated wave speeds satisfy $S_L = -c_L \le 0 \le c_R = S_R$ and the middle speed $S_\star$ works out to exactly $0$, which collapses the star-region states to something that leaves only the pressure term surviving in the momentum flux, with mass and energy fluxes vanishing identically. This ``nothing is happening, only pressure pushes on the face'' behavior is precisely the discrete analogue of hydrostatic balance, and it's the property the well-balanced modification is built to exploit.

The fix: instead of feeding HLLC the raw conserved states, we rescale each side by the ratio of a face-averaged equilibrium pressure to that side's own cell-interior equilibrium pressure:
\[
\widehat{\mathbf{F}} = \mathbf{F}^{hllc}\left(\frac{p_h^{e,\star}}{p_h^{e,\mathrm{int}(K)}}\,\mathbf{U}_h^{\mathrm{int}(K)},\ \ \frac{p_h^{e,\star}}{p_h^{e,\mathrm{ext}(K)}}\,\mathbf{U}_h^{\mathrm{ext}(K)}\ ;\ \mathbf{n}_{\mathcal{E},K}\right), \qquad p_h^{e,\star} = \tfrac12\left(p_h^{e,\mathrm{int}(K)} + p_h^{e,\mathrm{ext}(K)}\right).
\]
Why this particular rescaling? Because at exact equilibrium, $\mathbf{U}_h^{\mathrm{int}(K)}$ and $\mathbf{U}_h^{\mathrm{ext}(K)}$ are both zero-velocity states with pressures $p_h^{e,\mathrm{int}(K)}$ and $p_h^{e,\mathrm{ext}(K)}$ respectively, and generally $p_h^{e,\mathrm{int}(K)} \ne p_h^{e,\mathrm{ext}(K)}$ — the equilibrium pressure genuinely varies across a face, that's just what a stratified star looks like. The contact property only fires cleanly when both sides share the \emph{same} pressure. So we rescale each side up (or down) to the common reference value $p_h^{e,\star}$ before calling HLLC, apply the contact property there (where it gives exactly $(0,\,p_h^{e,\star}\mathbf{n}^T,\,0)^T$), and this becomes the discrete flux we actually use. The rescaling is undone implicitly, because it's baked directly into the momentum source term derivation next, so the two exactly cancel at equilibrium.

\paragraph{The well-balanced source term.} On the momentum equation, the full source is $\mathbf{S}^{[2]} = -\rho\grad\phi = -\rho\grad\phi^e - \rho\grad\phi^\delta$. The perturbation piece is discretized in the ordinary, direct way. The equilibrium piece is where the actual balancing act happens, and it's worth walking through the algebra in full, because every line of it exists for a specific reason:
\[
\int_K -\rho\grad\phi^e\, v\, d\mathbf{x}
= \int_K \rho\,\frac{\grad p^e}{\rho^e}\, v\, d\mathbf{x}
\]
(using the exact equilibrium relation $\grad p^e = -\rho^e\grad\phi^e$ to trade $\grad\phi^e$ for $\grad p^e/\rho^e$, a quantity we can evaluate directly from the known equilibrium profile). Now add and subtract the cell-average ratio $\overline{\rho}_K/\overline{\rho^e}_K$ inside the integral:
\[
= \int_K \grad p^e\left(\frac{\rho}{\rho^e} - \frac{\overline{\rho}_K}{\overline{\rho^e}_K}\right) v\, d\mathbf{x} \ +\ \frac{\overline{\rho}_K}{\overline{\rho^e}_K}\int_K \grad p^e\, v\, d\mathbf{x}.
\]
The first term is the genuinely new part of the perturbation source, and notice it vanishes identically the instant $\rho/\rho^e$ is constant across the cell — exactly what happens at equilibrium, where $\rho_h = \rho_h^e$ everywhere. The second term is integrated by parts, using precisely the same weak-form move that built DG in the first place, turning a volume integral of a gradient into a boundary flux minus a volume integral against $\grad v$:
\[
\frac{\overline{\rho}_K}{\overline{\rho^e}_K}\int_K \grad p^e\, v\, d\mathbf{x} = \frac{\overline{\rho}_K}{\overline{\rho^e}_K}\left(\sum_{\mathcal{E}\in\partial K}\int_{\mathcal{E}} p^e\, v\, \mathbf{n}_{\mathcal{E},K}\, ds - \int_K p^e\, \grad v\, d\mathbf{x}\right).
\]
Putting it together, and writing $\langle \mathbf{S}^{[2]},v\rangle_K$ for the fully quadrature-discretized version (quadrature points $\mathbf{x}_K^{(q)}$ with weights $w_K^{(q)}$ inside $K$, and $\mathbf{x}_\mathcal{E}^{(\mu)}$ with weights $w_\mathcal{E}^{(\mu)}$ on each face):
\[
\langle\mathbf{S}^{[2]},v\rangle_K =
\sum_{q} w_K^{(q)}\left(\frac{\rho_h(\mathbf{x}_K^{(q)})}{\rho_h^e(\mathbf{x}_K^{(q)})} - \frac{\overline{\rho}_K}{\overline{\rho^e}_K}\right)\grad p_h^e(\mathbf{x}_K^{(q)})\, v(\mathbf{x}_K^{(q)})
\]
\[
+\ \frac{\overline{\rho}_K}{\overline{\rho^e}_K}\left(\sum_{\mathcal{E}\in\partial K}\sum_{\mu} w_\mathcal{E}^{(\mu)} p_h^e(\mathbf{x}_\mathcal{E}^{(\mu)})\, v^{\mathrm{int}(K)}\, \mathbf{n}_{\mathcal{E},K} - \sum_{q} w_K^{(q)} p_h^e(\mathbf{x}_K^{(q)})\, \grad v(\mathbf{x}_K^{(q)})\right)
-\sum_{q} w_K^{(q)}\, \rho_h(\mathbf{x}_K^{(q)})\, \mathbf{g}_h^\delta(\mathbf{x}_K^{(q)})\, v(\mathbf{x}_K^{(q)}).
\]
That last line is just the ordinary perturbation-gravity contribution, $-\rho\grad\phi^\delta$, tacked back on with no special treatment. The whole point of everything above it is that, at exact equilibrium, this discrete source term is built term-for-term to exactly cancel the modified HLLC flux from the previous paragraph. That's not a coincidence or a lucky numerical accident; it's a designed algebraic identity, and it's the reason we can actually \emph{prove} the well-balanced property later instead of just observing it empirically on a few test cases.

\subsubsection{Fix two: a conservative form for total energy}

For energy conservation, the move is to stop evolving the ordinary non-gravitational energy $E$ directly, and instead evolve a total energy that folds in gravitational potential energy from the start:
\[
E_{tot} = E + \tfrac12 \rho\phi.
\]
Why does this help? Because with some care you can rewrite the evolution equation for $E_{tot}$ in genuinely conservative form — divergence of a flux, with no leftover volumetric source term at all:
\[
\pdv{E_{tot}}{t} + \grad\cdot\Big((E+p)\mathbf{u} + \mathbf{F}_g\Big) = 0,
\]
where the new piece, the \emph{gravity energy flux}, is
\[
\mathbf{F}_g = \frac{1}{8\pi G}\Big(\phi\,\grad\dot\phi - \dot\phi\,\grad\phi\Big) + \rho\mathbf{u}\phi.
\]
It's worth actually seeing why this identity is true, rather than taking it on faith, because the derivation is short and it explains where every piece of $\mathbf{F}_g$ comes from. Start from the time derivative of the gravitational energy piece, $\partial_t(\tfrac12\rho\phi)$, and add and subtract the divergence of the proposed flux and the ordinary source term $-\rho\mathbf{u}\cdot\grad\phi$ from the plain energy equation:
\[
\pdv{}{t}\left(\tfrac12\rho\phi\right) + \grad\cdot\mathbf{F}_g - \rho\mathbf{u}\cdot\grad\phi
= \tfrac12\dot\rho\,\phi + \tfrac12\dot\phi\,\rho + \frac{1}{8\pi G}\Big(\phi\,\Delta\dot\phi - \dot\phi\,\Delta\phi\Big) + \grad\cdot(\rho\mathbf{u})\,\phi.
\]
Now use mass conservation, $\dot\rho = -\grad\cdot(\rho\mathbf{u})$, to cancel the $\phi\,\dot\rho$ term against the last term above, and use $\Delta\phi = 4\pi G\rho$ together with its time-derivative $\Delta\dot\phi = 4\pi G\dot\rho = -4\pi G\,\grad\cdot(\rho\mathbf{u})$ to rewrite the remaining two terms. Every surviving term cancels in pairs, and the whole expression collapses to exactly zero. Add this identity to the ordinary energy equation and you land on the conservative form above. It's a genuinely satisfying piece of algebra: the messy, seemingly ad hoc combination of terms in $\mathbf{F}_g$ isn't arbitrary at all, it's precisely what's needed to make two independent physical statements (mass conservation and Poisson's equation) conspire to produce a clean conservation law for energy.

There's an obvious wrinkle: $\mathbf{F}_g$ contains $\dot\phi = \partial\phi/\partial t$, a time derivative of a field, which isn't something we normally track explicitly in a DG code. The way around it is to note that $\dot\phi$ itself satisfies another Poisson equation, obtained by differentiating $\Delta\phi = 4\pi G\rho$ in time and using mass conservation again:
\[
\Delta\dot\phi = 4\pi G\,\dot\rho = -4\pi G\,\grad\cdot(\rho\mathbf{u}).
\]
So we solve one additional elliptic problem, every stage, using the exact same LDG mixed formulation as before, now with $-4\pi G\,\grad\cdot(\rho\mathbf{u})_h$ as the source instead of $4\pi G\rho_h$, to obtain both $\dot\phi_h$ and its gradient $\dot{\mathbf{g}}_h$.

There's one more, purely numerical, subtlety worth being upfront about, because it's the kind of thing that will cost you a full order of accuracy if you don't know to watch for it. If you discretize the right-hand side $-4\pi G\,\grad\cdot(\rho\mathbf{u})_h$ naively — literally differentiating the DG polynomial $(\rho\mathbf{u})_h$ pointwise inside each cell — you lose an order of accuracy, because differentiating a degree-$k$ polynomial approximation generically only gives you a degree-$(k-1)$-accurate approximation of the true derivative. The fix, again, is to reach for the DG weak form itself rather than a pointwise derivative: apply summation-by-parts to move that derivative off of $(\rho\mathbf{u})_h$ and onto the test function $\psi$ instead, at the cost of picking up a boundary term built from the already-computed mass flux $\widehat{f}^{[1]}$:
\[
\int_K \dot{\mathbf{g}}_h\cdot\grad\psi\, d\mathbf{x} = \int_{\mathcal{E}} \psi^{\mathrm{int}(K)}\,\widehat{\dot{\mathbf{g}}}_h\cdot\mathbf{n}_{\mathcal{E},K}\, ds + \int_K 4\pi G\,(\rho\mathbf{u})_h\cdot\grad\psi\, d\mathbf{x} - \sum_{\mathcal{E}\in\partial K}\int_{\mathcal{E}} 4\pi G\, \widehat{f}^{[1]}_{\mathbf{n}_{\mathcal{E},K}}\, \psi^{\mathrm{int}(K)}\, ds.
\]
This one small technical patch restores the full design order of accuracy for $\dot\phi_h$, and by extension for the whole energy flux; skip it and you'll still get a scheme that conserves energy (that property doesn't depend on this fix), but one whose overall convergence rate quietly degrades, especially on coarser meshes.

With $E_{tot}$ as the evolved variable, we can repackage the whole system uniformly:
\[
\mathbf{W} = \begin{bmatrix}\rho\\ \rho\mathbf{u}\\ E_{tot}\end{bmatrix}, \quad
\mathbf{F}(\mathbf{W}) = \begin{bmatrix}\rho\mathbf{u}\\ \rho\mathbf{u}\otimes\mathbf{u}+p\mathbf{I}\\ (E+p)\mathbf{u}\end{bmatrix}, \quad
\mathbf{G}(\mathbf{W}) = \begin{bmatrix}\mathbf{0}\\ \mathbf{0}\\ \mathbf{F}_g\end{bmatrix}, \quad
\mathbf{S} = \begin{bmatrix}0\\ -\rho\grad\phi\\ 0\end{bmatrix},
\]
\[
\pdv{\mathbf{W}}{t} + \grad\cdot\mathbf{F}(\mathbf{W}) + \grad\cdot\mathbf{G}(\mathbf{W}) = \mathbf{S}.
\]

One bookkeeping detail that turns out to matter enormously in practice: at every stage we still need the ordinary energy $E$ (to get pressure out of the equation of state), but we're now evolving $E_{tot}$. The conversion is $E_h = (E_{tot})_h - \mathbf{P}\big(\tfrac12\rho_h\phi_h\big)$ — and that $\mathbf{P}$, the $L^2$ projection, is not optional decoration. If you instead compute $\tfrac12\rho_h\phi_h$ pointwise at quadrature points and skip projecting the product back into the DG polynomial space $V_h^k$, the formula \emph{looks} identical but the well-balanced cancellation from the previous section silently breaks, because the pressure recovered this way is no longer guaranteed to equal $p_h^e$ exactly at equilibrium. This is exactly the kind of easy-to-miss implementation detail that separates ``it looks right and mostly works'' from ``it's provably well-balanced,'' and it's the sort of thing I'd genuinely only have caught by trying to prove the property myself and watching the proof break.

\subsection{Time Discretization}

Everything above was written semi-discretely: space is discretized into the DG polynomial spaces, but time is still continuous, giving us an ODE system of the schematic form
\[
\mathbf{W}_h^{n+1} = \mathbf{W}_h^n + \Delta t\, \mathcal{H}(\mathbf{W}_h^n, t^n),
\]
where $\mathcal{H}$ is shorthand for the entire right-hand side we just spent two sections deriving: the modified HLLC flux, the well-balanced source term, and the gravity energy flux. (The standard, non-structure-preserving scheme has an analogous operator, which I'll call $\mathcal{L}(\mathbf{U}_h^n,t^n)$.) Forward Euler, as written, is only first order accurate in time, which would badly undercut all the care we just took getting high spatial order — there's no point running a $P^3$ DG method in space if a first-order time integrator is going to wash out the accuracy anyway. The standard fix, and the one used throughout this scheme, is strong-stability-preserving Runge-Kutta (SSP-RK), specifically chosen because it's built as a convex combination of forward-Euler steps, which means any stability or positivity property you can prove for forward Euler carries over automatically to the full multistage method, with no extra work.

Second order SSP-RK:
\[
\mathbf{W}_h^{(1)} = \mathbf{W}_h^n + \Delta t\,\mathcal{H}(\mathbf{W}_h^n, t^n), \qquad
\mathbf{W}_h^{n+1} = \tfrac{1}{2}\mathbf{W}_h^n + \tfrac{1}{2}\Big(\mathbf{W}_h^{(1)} + \Delta t\,\mathcal{H}(\mathbf{W}_h^{(1)}, t^n+\Delta t)\Big).
\]

Third order SSP-RK:
\[
\mathbf{W}_h^{(1)} = \mathbf{W}_h^n + \Delta t\,\mathcal{H}(\mathbf{W}_h^n, t^n),
\]
\[
\mathbf{W}_h^{(2)} = \tfrac34 \mathbf{W}_h^n + \tfrac14\Big(\mathbf{W}_h^{(1)} + \Delta t\,\mathcal{H}(\mathbf{W}_h^{(1)}, t^n+\Delta t)\Big),
\]
\[
\mathbf{W}_h^{n+1} = \tfrac13 \mathbf{W}_h^n + \tfrac23\Big(\mathbf{W}_h^{(2)} + \Delta t\,\mathcal{H}(\mathbf{W}_h^{(2)}, t^n+\tfrac12\Delta t)\Big).
\]

Each stage is literally just another forward-Euler-style evaluation of $\mathcal{H}$, meaning the entire elliptic solve, modified HLLC flux, and gravity-energy-flux pipeline from the previous section gets called again, unchanged, at every stage. This reusability is one of the more pleasant features of how the scheme was built: none of the well-balanced or energy-conserving modifications care what order of Runge-Kutta method is wrapped around them, so raising the temporal order later is purely a matter of adding stages, with zero changes to the spatial operator itself.

Practically, the timestep is chosen from a standard CFL condition based on the fastest wave speed present in any cell,
\[
\Delta t = \mathrm{CFL} \cdot \frac{h}{\max_K\left(\abs{\overline{\mathbf{u}}_K}_\infty + \overline{c}_K\right)}, \qquad \overline{c}_K = \sqrt{\gamma\, \overline{p}_K/\overline{\rho}_K},
\]
with $h$ the smallest element diameter in the mesh and $\overline{(\cdot)}_K$ denoting the cell average. Notice gravity doesn't directly enter this CFL restriction — the elliptic solve is implicit (we're solving a linear system exactly, not marching it forward explicitly), so the only explicit wave-speed restriction left is the ordinary hydrodynamic one.

\subsection{Ossilation Elimination}

High order polynomials are wonderful at representing smooth flow and genuinely terrible at representing a shock. Near a real discontinuity — the blast wave from our earlier explosion example, say — a high-degree DG polynomial will ring: classic Gibbs-phenomenon over- and under-shoots that, in the worst case, can drive density or pressure negative at a quadrature point and crash the whole simulation outright.

The traditional fix here is a slope limiter, but limiters can be blunt instruments: applied indiscriminately, they can easily destroy the very fine cancellations we spent the last section engineering into the well-balanced source term and modified flux. Instead, this scheme reaches for a more recent idea, the oscillation-eliminating (OE) technique, which works by letting the high-degree moments of the polynomial solution decay through an artificial damping ODE, applied as a separate fractional step tacked on after each Runge-Kutta stage, rather than by clipping or reconstructing the polynomial directly
\[
\dv{\mathbf{U}}{t} + \mathbf{L}_f(\mathbf{U}) + \mathbf{\Sigma}(\mathbf{U})\mathbf{U} = 0
\quad\Longrightarrow\quad
\dv{\mathbf{U}}{t} = -\mathbf{L}_f(\mathbf{U}), \qquad \dv{\mathbf{U}_\sigma}{t} = -\mathbf{\Sigma}(\mathbf{U})\mathbf{U}_\sigma.
\]
The first piece is just the ordinary DG spatial operator we already have. The second, $\mathbf{\Sigma}(\mathbf{U})\mathbf{U}_\sigma$, is a purely local, per-cell damping term, and its coefficient $\vartheta_K^m$ is built from a smoothness indicator, summed over each cell's faces:
\[
\vartheta_K^m(\mathbf{U}_h) = \sum_{e\in\partial K} \beta_e\, \frac{\sigma_{e,K}^m(\mathbf{U}_h)}{h_{e,K}},
\]
where $\beta_e$ is a local estimate of the maximum wave speed across face $e$ (the spectral radius of the flux Jacobian in the normal direction), $h_{e,K}$ is roughly the distance from face $e$ to the far side of the cell, and the indicator itself,
%\[
%\sigma_{e,K}^m(u_h) = \frac{(2m+1)h_{e,K}^m}{2(2k-1)\,m!}\sum_{|\alpha|=m}\frac{\frac{1}{|e|}\int_e \left|[\partial^\alpha u_h]_e\right|\, dS}{\|u_h-\mathrm{avg}(u_h)\|_{L^\infty(\Omega)}}
%\]
measures how big the jump in the $m$-th derivative across that face is, relative to how much the solution varies overall. In a smooth region, that jump is tiny compared to the solution's global variation, so $\vartheta_K^m$ stays small and the damping barely touches anything. Right next to a shock, the jump is large, $\vartheta_K^m$ spikes, and the damping aggressively suppresses the offending high-order moments — but, and this is the important structural feature, it leaves the lowest moment, the cell average, completely untouched by construction ($\vartheta_K^0$ effectively doesn't act on the mean). The damping equation actually has a closed form: writing the solution in a hierarchical orthogonal basis $\phi_K^{(\alpha)}$, each higher moment simply decays exponentially at a rate set by the local indicator,
\[
\mathcal{F}_{\Delta t}\mathbf{U}_h = \mathbf{U}_K^{(0)}\phi_K^{(0)}(\mathbf{x}) + \sum_{j=1}^k e^{-\Delta t \sum_{m=0}^j \vartheta_K^m(\mathbf{U}_h)} \sum_{|\alpha|=j} \mathbf{U}_K^{(\alpha)}\phi_K^{(\alpha)}(\mathbf{x}).
\]
Because the cell average $\mathbf{U}_K^{(0)}$ is left alone, mass, momentum, and total energy on each cell are completely unaffected by the OE step; only the shape of the polynomial \emph{within} the cell gets smoothed. That's exactly the property we need to keep conservation intact.

The remaining piece of care, specific to this well-balanced scheme, is where the damping is allowed to act. It's applied only to the perturbation part $\mathbf{W}^\delta$ of the solution, never to the frozen equilibrium background $\mathbf{W}^e$:
\[
\mathbf{W}_h = \mathbf{W}_h^e + \mathcal{F}_{\Delta t}(\mathbf{W}_h^\delta).
\]
If the damping were instead applied to the full solution, including the equilibrium piece, it would gradually smooth away the very steady-state structure the well-balanced property exists to protect — a hydrostatic star has real, physically meaningful higher moments in its density profile, and those aren't the ``oscillations'' we're trying to eliminate. Restricting the damping to the perturbation alone means: on a problem that's exactly at equilibrium, $\mathbf{W}^\delta \equiv 0$, the damping has literally nothing to act on, and the well-balanced property survives the OE step trivially. On a genuinely perturbed or shocked problem, one can show the damping only introduces an error of the same order as the underlying DG truncation error itself, so smooth problems don't quietly lose accuracy — you only pay a price near an actual discontinuity, which is exactly where you want to be paying it.

For the very strongest shocks — the blast wave test later, where the central pressure jumps by two orders of magnitude — OE alone isn't quite enough to guarantee density and pressure stay positive everywhere, so it's paired with a separate, more surgical positivity-preserving (PP) limiter, which only activates at a quadrature point if density or pressure there is actually about to go negative, and otherwise leaves the solution untouched.

\subsection{Structure Preserving Properties}

It's worth stating plainly, and proving carefully, exactly what all of this extra machinery buys us, because the algebra above is dense enough that the punchline can get lost.

\subsubsection{Theorem: well-balancedness}

\emph{Claim.} If the scheme is initialized exactly at a known hydrostatic equilibrium — $\mathbf{u}^0 = \mathbf{0}$, $p^0 = -\rho^0\grad\phi^0$, $\Delta\phi^0 = 4\pi G\rho^0$ — then for every subsequent timestep $n$,
\[
\rho_h^n = \rho_h^e, \qquad (\rho\mathbf{u})_h^n = \mathbf{0}, \qquad p_h^n = p_h^e, \qquad \phi_h^n = \phi_h^e,
\]
exactly, up to floating-point round-off, regardless of mesh resolution.

\emph{Proof sketch, by induction on $n$.} Assume the claim holds at step $n$. Walk through the flowchart from the next section term by term. Since $\rho_h^n = \rho_h^e$, the perturbation density is exactly zero, so the perturbation Poisson solve returns $\phi_h^\delta = 0$ and hence $\phi_h^n = \phi_h^e$ exactly. From this, the recovered non-gravitational energy is $E_h^n = p_h^e/(\gamma-1)$ and the recovered pressure is exactly $p_h^e$ — this is precisely where the consistent $L^2$-projection bookkeeping flagged earlier becomes essential; without it this line fails. Feeding $\mathbf{U}_h^{n,\mathrm{int}(K)} = (\rho_h^{e,\mathrm{int}(K)},\mathbf{0},p_h^{e,\mathrm{int}(K)}/(\gamma-1))^T$ (and the analogous exterior state) into the modified HLLC flux, the contact property fires and gives $\widehat{\mathbf{F}}^n = (0,\, p_h^{e,\star}\mathbf{n}^T,\, 0)^T$ — in particular the mass-flux component $\widehat{f}^{n,[1]}=0$. Since the mass flux is zero everywhere, the second Poisson solve for $\dot\phi_h,\dot{\mathbf{g}}_h$ has a zero right-hand side, so $\dot\phi_h^n = 0$, $\dot{\mathbf{g}}_h^n = \mathbf{0}$, and the entire gravity energy flux vanishes identically, $(\widehat{\mathbf{F}_g})^n = 0$. Finally, substitute all of this into the momentum update: the four resulting pieces — the pressure-gradient volume term, the ``add-and-subtract'' perturbation term from the source, the rescaled-pressure boundary term, and the perturbation-gravity term — cancel exactly in pairs, term for term, leaving $(\rho\mathbf{u})_h^{n+1} = (\rho\mathbf{u})_h^n = \mathbf{0}$. Since the mass and energy fluxes were already shown to vanish, $\rho_h^{n+1}=\rho_h^e$ and $(E_{tot})_h^{n+1} = (E_{tot})_h^n$ follow immediately, closing the induction. $\blacksquare$

The proof is mechanical once you have it laid out, but the punchline is genuinely strong: this isn't an asymptotic statement about the error shrinking as $h\to0$. The equilibrium is preserved exactly, on any mesh, at any polynomial order, because every single term that could disturb it has been shown to vanish identically rather than merely become small.

\subsubsection{Theorem: total energy conservation}

\emph{Claim.} Under periodic boundary conditions, or compactly-supported boundary conditions (velocity vanishing on the boundary), the scheme conserves discrete total energy exactly at every timestep:
\[
\int_\Omega (E_{tot})_h^{n+1}\, d\mathbf{x} = \int_\Omega (E_{tot})_h^n\, d\mathbf{x}.
\]

\emph{Proof sketch.} Take the discrete $E_{tot}$ update equation, set the test function $v=1$, and sum over every cell $K$ in the mesh. Every interior face appears in this sum exactly twice, once from each of the two cells sharing it, with opposite orientations of the outward normal — so every interior-face contribution from the energy flux and the gravity energy flux cancels exactly against its neighbor: whatever leaves one cell enters the other, unchanged. What's left after that cancellation is purely a sum over boundary faces $\mathcal{E}_B$. Under periodic boundary conditions, boundary faces are identified in matching pairs on opposite sides of the domain and cancel each other the same way interior faces do. Under compactly-supported boundary conditions, velocity vanishes identically on $\mathcal{E}_B$, which makes both the ordinary energy flux and the gravity energy flux vanish there directly, since every term in each flux carries at least one explicit factor of $\mathbf{u}$. Either way, the total boundary contribution is exactly zero, and we're left with $\int_\Omega (E_{tot})_h^{n+1} = \int_\Omega (E_{tot})_h^n$. $\blacksquare$

Notice this proof didn't need the well-balanced machinery at all — TEC is a genuinely independent property, coming entirely from writing the energy equation in conservative form. In practice this shows up as a total-energy trace that's flat to machine precision on a well-resolved, stable configuration, and — arguably more usefully — as an energy diagnostic you can actually trust while something dramatic is unfolding, like a gravitational collapse, precisely the regime where the standard scheme's energy error tends to grow fastest and become least trustworthy.

\subsection{The Code itself (Flowchart)}

Putting all of the above in the order you'd actually implement it, here is one full Runge-Kutta stage, advancing from a known state $\big(\rho_h^n, (\rho\mathbf{u})_h^n, (E_{tot})_h^n\big)$ to the corresponding stage value:

\begin{itemize}
  \item \textbf{Step 0 --- initialization (once, at $t=0$ only).} Project the physical initial conditions into the DG space: $\rho_h^0 = \mathbf{P}\rho^0$, $(\rho\mathbf{u})_h^0 = \mathbf{P}(\rho\mathbf{u})^0$, $E_h^0 = \mathbf{P}E^0$. Solve the equilibrium Poisson problem once, $(\mathbf{g}_h^e,\phi_h^e) = \mathcal{D}_1(4\pi G\rho_h^e)$, and freeze it for the rest of the run. Solve the initial perturbation Poisson problem to get $\phi_h^0$, and set $(E_{tot})_h^0 = E_h^0 + \mathbf{P}\big(\tfrac12\rho_h^0\phi_h^0\big)$ — using the projection here, not a pointwise product, for exactly the reason flagged above.
  \item \textbf{Step 1 --- gravity.} Given $\rho_h^n$, solve the perturbation Poisson problem $(\mathbf{g}_h^\delta, \phi_h^\delta)^n = \mathcal{D}_1\big(4\pi G(\rho_h^n - \rho_h^e)\big)$ and recombine with the frozen equilibrium piece to get the total field, $(\mathbf{g}_h^n,\phi_h^n) = (\mathbf{g}_h^e,\phi_h^e) + (\mathbf{g}_h^\delta,\phi_h^\delta)^n$. Separately, solve the summation-by-parts-corrected Poisson problem $\mathcal{D}_2$ for the time-derivative fields, $(\dot{\mathbf{g}}_h^n,\dot\phi_h^n) = \mathcal{D}_2\big(4\pi G\,\grad\cdot(\rho\mathbf{u})_h^n\big)$. Recover the ordinary energy, $E_h^n = (E_{tot})_h^n - \mathbf{P}\big(\tfrac12\rho_h^n\phi_h^n\big)$, and from it the pressure via the equation of state.
  \item \textbf{Step 2 --- fluid flux.} At every interior face, compute the equilibrium-pressure-rescaled states on each side and evaluate the modified HLLC flux $\widehat{\mathbf{F}}^n$ from those rescaled states.
  \item \textbf{Step 3 --- gravity energy flux.} Using $\phi_h^n,\mathbf{g}_h^n$ from Step 1, $\dot\phi_h^n,\dot{\mathbf{g}}_h^n$ also from Step 1, and the mass-flux component $\widehat{f}^{n,[1]}$ from Step 2, assemble both the volume contribution to $\mathbf{F}_g^n$ inside each element and its face contribution $(\widehat{\mathbf{F}_g})^n$.
  \item \textbf{Step 4 --- assemble and update.} Combine the fluid flux (Step 2), the well-balanced momentum source term (built the same way as in Step 1's rescaling), and the gravity energy flux (Step 3) into the full right-hand side $\mathcal{H}$, and take one Runge-Kutta stage update for $\big(\rho_h,(\rho\mathbf{u})_h,(E_{tot})_h\big)$.
  \item \textbf{Step 5 --- oscillation control.} Apply the OE damping (and, for strongly discontinuous problems, the positivity-preserving limiter) to the perturbation part of the newly updated solution only, leaving the frozen equilibrium piece untouched.
\end{itemize}

Repeat Steps 1 through 5 once per Runge-Kutta stage (twice for the second-order method, three times for third order), then move on to the next timestep with a freshly computed CFL-limited $\Delta t$. The only genuinely expensive part of this whole loop is the pair of elliptic solves in Step 1 — but because the mesh is fixed for the entire simulation, the factorization of the underlying stiffness matrix (Cholesky for the compactly-supported case, LU for the periodic case) is computed exactly once, before the time loop even starts, and reused at every stage of every timestep, which keeps the per-step cost dominated by the same sparse triangular solves you'd need for any implicit elliptic problem, not by repeated factorization.

\subsection{Results}

Here I have the results in the form of fun cool animations
%%%%%%%%%%%%%%%%%%%%%%%%%%%%%%%%%%%%%%%
\section{What is Velocity(joke video)}
\begin{itemize}
        \item Basic
  \item Average speed
  \item Derivative form
  \item Vector form
  \item Newtonian
  \item Newtonian + inertial refrence frame
  \item Newtonian Covariance/Galilean Relativity
  \item Newtonian-Cartan/ general covariance
  \item Fluid
  \item Special Relativity
  \begin{itemize}
    \item Coordinate Velocity
    \item Four-Velocity
    \item Proper Velocity
    \item Rapidity
    \item Relative Velocity
  \end{itemize}
  \item General Relativity
  \begin{itemize}
    \item Coordinate velocity
    \item GR four-velocity
    \item Tangent vector
    \item Covariant velocity
    \item Geodesic
    \item Warp Drive
  \end{itemize}
  \item Dynamics
  \item Lagrandian
  \item Hamiltonian
  \item Phase-Space Flow
  \item Poisson Structure
  \item Symplectic flow
\end{itemize}


\subsection{Introduction}

Let's define velocity, I mean its a pretty simple concept, you could ask a fifth grader what it is and they could define it with ease. I mean it is the simple equation...
\[
\begin{aligned}
V
&\equiv
(Q_{\xi},X)
=
\mathfrak{X}_{Q_{\xi}}[X]
\\[1ex]
\delta[Q_{\xi}]
&=
K(\delta,\mathfrak{X}_{Q_{\xi}})
\\
&=
\Omega_{\mathrm{ext}}(\delta,\mathfrak{X}_{Q_{\xi}})
-
G(\delta,\mathfrak{X}_{Q_{\xi}})
\\[2ex]
&=
\Omega(\delta,\mathfrak{L}_{\xi})
+
Q_{[\mathbb{X}(\delta),\mathbb{X}(\mathfrak{L}_{\xi})]}
+
\oint_{\partial\Sigma}
\mathbb{X}(\mathfrak{L}_{\xi})
\mathbin{\lrcorner}
\vartheta(\delta)
\\
&\qquad
+
\oint_{\partial\Sigma}
\mathbb{X}(\mathfrak{L}_{\xi})
\mathbin{\lrcorner}
\vartheta(\delta)
-
G(\delta,\mathfrak{L}_{\xi})
\\[2ex]
&=
\Omega(\delta,\mathfrak{L}_{\xi})
+
Q_{[\mathbb{X}(\delta),\mathbb{X}(\mathfrak{L}_{\xi})]}
+
\oint_{\partial\Sigma}
\mathbb{X}(\mathfrak{L}_{\xi})
\mathbin{\lrcorner}
\vartheta(\delta)
\\
&\qquad
+
\oint_{\partial\Sigma}
\mathbb{X}(\mathfrak{L}_{\xi})
\mathbin{\lrcorner}
\vartheta(\delta)
\\[2ex]
&=
\Omega_{\mathrm{ext}}(\delta,\mathfrak{L}_{\xi})
-
G(\delta,\mathfrak{L}_{\xi})
\\[2ex]
&=
\Omega(\delta,\mathfrak{L}_{\xi})
+
Q_{[\mathbb{X}(\delta),\mathbb{X}(\mathfrak{L}_{\xi})]}
+
\oint_{\partial\Sigma}
\mathbb{X}(\delta)
\mathbin{\lrcorner}
\vartheta(\mathfrak{L}_{\xi})
\\
&\qquad
-
\Omega(\mathfrak{L}_{\xi},\delta)
-
Q_{[\mathbb{X}(\mathfrak{L}_{\xi}),\mathbb{X}(\delta)]}
-
\oint_{\partial\Sigma}
\mathbb{X}(\mathfrak{L}_{\xi})
\mathbin{\lrcorner}
\vartheta(\delta)
\\
&\qquad
-
\left[
-\oint_{\partial\Sigma}
\mathbb{X}(\delta)
\mathbin{\lrcorner}
\vartheta(\mathfrak{L}_{\xi})
-
\oint_{\partial\Sigma}
\mathbb{X}(\mathfrak{L}_{\xi})
\mathbin{\lrcorner}
\vartheta(\delta)
\right]
\\[2ex]
\end{aligned}
\]
\[
\begin{aligned}
&=
\delta
\left[
\int_{\Sigma}
\left(
\vartheta(\mathfrak{L}_{\xi})
-
\xi\mathbin{\lrcorner}L
\right)
\right]
-
Q_{\delta[\xi]}
-
\oint_{\partial\Sigma}
\xi\mathbin{\lrcorner}\vartheta(\delta)
\\
&\qquad
+
Q_{[\mathbb{X}(\delta),\mathbb{X}(\mathfrak{L}_{\xi})]}
+
\oint_{\partial\Sigma}
\mathbb{X}(\mathfrak{L}_{\xi})
\mathbin{\lrcorner}
\vartheta(\delta)
\\[2ex]
&=
\delta
\left[
\int_{\Sigma}
\left(
\vartheta(
\mathcal{L}_{\xi}g_{ab},
\mathcal{L}_{\xi}\psi^{I}
)
-
\xi\mathbin{\lrcorner}L
\right)
\right]
-
Q_{\delta[\xi]}
\\
&\qquad
+
Q_{[\mathbb{X}(\delta),\mathbb{X}(\mathfrak{L}_{\xi})]}
+
\oint_{\partial\Sigma}
\mathbb{X}(\mathfrak{L}_{\xi})
\mathbin{\lrcorner}
\vartheta(\delta)
\\[2ex]
&=
\delta
\left[
\int_{\Sigma}
\left(
\vartheta(
2\nabla_{(a}\xi_{b)},
\mathcal{L}_{\xi}\psi^{I}
)
-
\xi\mathbin{\lrcorner}L
\right)
\right]
-
Q_{\delta[\xi]}
\\
&\qquad
+
Q_{[\mathbb{X}(\delta),\xi]}
+
\oint_{\partial\Sigma}
\xi\mathbin{\lrcorner}\vartheta(\delta)
\\[2ex]
\end{aligned}
\]
\[
\begin{aligned}
&=
\delta
\left[
\int_{\Sigma}
\left(
\vartheta(
\nabla_{a}\xi_{b}
+
\nabla_{b}\xi_{a},
\mathcal{L}_{\xi}\psi^{I}
)
-
\xi\mathbin{\lrcorner}L
\right)
\right]
-
Q_{\delta[\xi]}
\\
&\qquad
+
Q_{[\mathbb{X}(\delta),\xi]}
+
\oint_{\partial\Sigma}
\xi\mathbin{\lrcorner}\vartheta(\delta)
\\[2ex]
&=
\delta
\left[
\int_{\Sigma}
\left(
\vartheta(
\partial_{a}\xi_{b}
+
\partial_{b}\xi_{a}
-
2\Gamma^{c}{}_{ab}\xi_{c},
\mathcal{L}_{\xi}\psi^{I}
)
-
\xi\mathbin{\lrcorner}L
\right)
\right]
\\
&\qquad
-
Q_{\delta[\xi]}
+
Q_{[\mathbb{X}(\delta),\xi]}
+
\oint_{\partial\Sigma}
\xi\mathbin{\lrcorner}\vartheta(\delta)
\\[2ex]
&=
\delta
\left[
\int_{\Sigma}
\left(
\vartheta(
\partial_{a}\xi_{b}
+
\partial_{b}\xi_{a}
-
g^{cd}
(
\partial_{a}g_{db}
+
\partial_{b}g_{da}
-
\partial_{d}g_{ab}
)
\xi_{c},
\mathcal{L}_{\xi}\psi^{I}
)
-
\xi\mathbin{\lrcorner}L
\right)
\right]
\\
&\qquad
-
Q_{\delta[\xi]}
+
Q_{[\mathbb{X}(\delta),\xi]}
+
\oint_{\partial\Sigma}
\xi\mathbin{\lrcorner}\vartheta(\delta)
\\[3ex]
\dot{F}
&=
(Q_{\xi},F)
\\[1ex]
(Q_{\xi},g_{ab})
&=
\mathcal{L}_{\xi}g_{ab}
=
2\nabla_{(a}\xi_{b)}
\\[1ex]
(Q_{\xi},\psi^{I})
&=
\mathcal{L}_{\xi}\psi^{I}
\\[1ex]
(Q_{\xi},x^{\mu})
&=
\xi^{\mu}
\\[2ex]
\Gamma^{c}{}_{ab}
&=
\frac{1}{2}
g^{cd}
(
\partial_{a}g_{db}
+
\partial_{b}g_{da}
-
\partial_{d}g_{ab}
)
\\[1ex]
g^{ac}g_{cb}
&=
\delta^{a}{}_{b},
\qquad
\nabla_{c}g_{ab}=0
\\[2ex]
(Q_{\xi},Q_{\xi})
&=
-G(\mathfrak{L}_{\xi},\mathfrak{L}_{\xi})
\\
&=
\oint_{\partial\Sigma}
\xi\mathbin{\lrcorner}
\vartheta(\mathfrak{L}_{\xi})
\end{aligned}
\]

... uhm, well... maybe I am wrong about that.

\subsection{Basic}

Ok, well then, let's start at the beginning. Velocity can generally be defined by the change in the distance traveled divided by the time it takes to travel.
\[
v= \frac{\Delta x}{\Delta t}
  \]
  If we take the infintesimal transformation of such we get
\[
v = \dv{x}{t}
\]

But, as your middle school science teacher likely drilled in your head, velocity isn't just speed, it also has a direction, making a vector the easiest way to represent it.
\[
\mathbf{\dot{r}} = \dv{\mathbf{r}}{t} = \langle \dv{x}{t}, \dv{y}{t}, \dv{z}{t} \rangle
  \]

\subsection{Newtonian}

Now that we have a basic understanding of the math behind velocity we can start putting into context. Likey eveyone here knows newtons second law
\[
\mathbf{F} = m \dv[2]{\mathbf{r}}{t}
  \]
  solving for velocity gives

\[
\dot{r} = \int \frac{\mathbf{F}}{m} dt
  \]

  Now this is all nice and good, but it forgets relativity, I don't mean special or general, I mean Galilean relativity. Keeping that in track is very simple

  \[
\dot{r}' = \int \frac{\mathbf{F}}{m} dt - \dot{r}_{\text{refrence frame}}
  \]

  This brings us to our next problem, as having Newtonian mechanics rely on frames of refrences comes up with several problems, I am sure many people have heard high school schience teachers telling you that there is no centrifugal forces, as it is a fictitious force created by frames of refrence, to solve this problem we make Newtonian mechanics general covariant by using the Newton-Cartan method. Essentially taking a space-time approach similar to general relativity but making time absolute rather than relative.

  Defining velocity as
\[
V^\mu = \dv{x^\mu}{dt} = (1, \dv{x^i}{dt})
  \]

  While this equation may look pretty simple, it is an extreme simplification. I really don't know the purpose of the Newton-Cartan other than to simply prove Einstien wrong.
\subsection{Special Relativity}

This talk or relativity branches cleanly into special relativity, as velocity truly begins to change. As time becomes another relative dimension to move through, which is a really funny idea to think about. Running through time.

Anyways, the velocity vector obviously becomes 4-d adding time to the mix

\[
U^\mu = \dv{x^\mu}{\tau}
  \]

  Where $\tau$ is proper time, or the subjective time of the travelers rather than those observing them. Take the tiny people on a relativistic baseball as compared to the batter waiting for them.

We can relatively easily then expanded this to our traditional relaitive speed.

Taking galelian relativity
\[
v_{\text{rel}} = v_1 - v_0
  \]

  and adding relativistic corrections
\[
v_{\text{rel}} = \frac{v_1 - v_0}{1 - v_1 v_2 / c^2}
  \]
Because why not
\subsection{GR}

But, this method isn't enough, there is now the problem of absolute acceleration, creating paradoxes and and preferential frames of refrence, something physically nonsensical in a relative world without absolute space or time. Though, using our earlier covariant newtonian framework and combining that with Einstienian proper time we get an equation that looks the same
\[
U^\mu = \dv{x^\mu}{\tau}
  \]

  but now
\[
g_{\mu \nu}U^\mu U^\nu = -c^2
  \]

  Thus adding a curved geometry of space-time that adds further separation between observers. The rest also stays the same as SR but with further corrections, and can be expressed through the language of remanian geometry rather than Lorentzian or geodesics, I won't get into all of that though.

  Though, before we move on, I want to talk about the velocity of warp drives. If you haven't heard, a solution of GR allows for apperent faster than speed of light movement, really messing with our definition of velocity which in all cases must be less than the speed of light.

  It doesn't really break anything as most solutions have it where the matter doesn't move at all but spacetime shrinks and grows behind it. Making its actual movement 0 m/s yet to an outside observer moving several times the speed of light as it moves a great ``distance'' in a very short time. Making definitions a bit confusing.

\subsection{Dynamics}

Ok, all of this spacetime stuff is really confusing and make a great deal of a mess. Lets move onto a completely different approach. Moving away from kinematics and into dynamics. Involving phase-space, where there aren't some defined set of dimensions of space and time, but an infinite amount of dimenisons of anything you could observe; momentum, position, time, heat, entropy, well... not really traditionally entropy. Anyways, rather than taking an infinite amount of dimensions, let's shrink it down to two, the two most generally used, momenta and position. Then we will add mass but as it is constant in our situation it is infinitely then and thus can be expressed in 2d

(show graph)

Here, momenta and position are not like their Newtonian cousians, but rather generalized. They are related pairs and aren't constrained in the way position and momentum generally are. Position can mean angle, position plus direction rolled in one, or all three dimensions rolled into one vector field. The same with momenta, but to make this example concret, we will take the traditional picture of position and momenta to give that if one represents the flow of the vector field as H, or commonly refered to as the Hamiltonain, we can express the derivative of position as a function of time as
\[
v = \dot{q} = \pdv{H}{p_i}
  \]

  Why, that is a really big question, just trust me.

  Thus saying that velocity is a part of the vector field on phasespace

  (show on graph)

  but what is phase space, well... that is an even bigger question but I will take a quick jab at it.

  First we introduce something called a Poisson Structure
\[
  \dot{f} = \pb{f}{H} = \pdv{f}{q^i} \pdv{H}{p_i} - \pdv{f}{p_i} \pdv{H}{q^i}
  \]

  Where the entire vector field can thus be written as
\[
X_H = \pb{\cdot}{H}
  \]
  Where the dot represents all possible observables, like position, momenta, or really anything else. The geometry behind the vector field is the phase-space. This geometry is something called symplectic, but I won't get into that.


\subsection{Other}

Now, let's speed run some others. Fluid velocity is a vector function that attaches a newtonian velocity to every point in the velocity.

Quantum mechanics adds fuzzy probability to velocity
\[
\mathbf{v} = \frac{\mathbf{j}}{\abs{\psi}^2}
  \]

  Kinetic theory combines probability of QM with the continuum of fluids

  Geometric mechanics defines it as a ``natural element of a tangent bundle,'' which seems hard to explain





\section{Developing Numerical General Relativity: Part 1, Vacuum Kerr Black Whole}

\subsection{Planning}
\begin{itemize}
  \item Introduction
  \item Scope
  \item Kerr SpaceTime
  \item 3+1 split
  \item ADM Formulism
  \item BSSN Formulism
  \item Initial data
  \item Cartoon Method
  \item Excision
  \item Method of lines + RK4
  \item Boundary conditions
  \item Code
  \item Validation
  \item Problems I had
  \item Next steps
  \item Animations
\end{itemize}

\subsection{Introduction}

Welcome to my first video on developing a numerical general relativistic code, starting with a vacuum Kerr Black Hole.

\subsection{Scope}

Before getting to the scope of this video, I will mention my overall scope of my broader project. I plan to create a FORTRAN code that uses General Relativity, Relativistic Fluid dynamics, and I am thinking about adding additional tools like magnetic or Nuclear EOS.

The main goal is for the final output to be a toy model of a quasi-star, a theoretical extremely large star with a black hole within its center but yet too large to undergo a supernova, this obviously combines a lot of extreme relativistic and fluid tools.

Rather than jumping straight into this extremely large project, I will do it in steps. I will start with a fairly basic one that uses just GR, a kerr black hole in a vacuum. That is what I am doing

\subsection{3+1 Split}

To solve most physics simulation, we derive them as initial value problems as a function of time and then numerically integrate. Those that have an understanding of GR can see how this is troubling.
